% Auto-extracted from ../tfpt-42.tex for the TFPT 4.5 clean split.
% Source body assigned by the TFPT 4.5 clean split.

% ---- Carrier algebraic normal form through abelian index: source lines 1306-2314 ----
\subsection{Carrier algebraic normal form}
\label{sec:minimal-carrier-proof}

\begin{lemma}[Algebraic normal form from exterior signatures]
\label{thm:carrier-minimality-k}
\label{thm:carrier-minimality-kpp}
Let
\[
E=E_-\oplus E_+
\]
be a finite chiral carrier with two simple factors of ranks
\[
b:=\dim E_-,
\qquad
s:=\dim E_+,
\]
and let
\[
S^+=\Lambda^{\mathrm{even}}E
\]
be the positive half-spinor packet. Let
\[
X=q_-P_-+q_+P_+
\]
be a primitive integer two-point generator with
\[
q_-<0<q_+,
\qquad
bq_-+sq_+=0.
\]
Assume only
\[
\dim \Lambda^2 E_+=1,
\qquad
\Lambda^2 E_- \cong E_-^\vee.
\]
Then necessarily
\[
s=2,
\qquad
b=3,
\qquad
g:=b+s=5.
\]
Hence
\[
\dim \Lambda^{\mathrm{even}}E=16.
\]
The even exterior algebra decomposes into exactly the six sectors
\[
(1,1)_0,\quad
(3,2)_{1/6},\quad
(\bar 3,1)_{-2/3},\quad
(\bar 3,1)_{1/3},\quad
(1,2)_{-1/2},\quad
(1,1)_1,
\]
and the primitive trace condition therefore has the unique solution
\[
(q_-,q_+)=(-2,3).
\]
Consequently
\[
X=-2P_3+3P_2,
\qquad
Y:=\frac{X}{6}
=
\mathrm{diag}\!\left(-\frac13,-\frac13,-\frac13,\frac12,\frac12\right).
\]
\end{lemma}

\begin{proof}
From
\[
\dim\Lambda^2 E_+=\binom{s}{2}=1
\]
one gets $s=2$. For the color factor one always has
\[
\Lambda^{b-1}E_- \cong E_-^\vee.
\]
The assumption
\[
\Lambda^2E_- \cong E_-^\vee
\]
therefore forces $b-1=2$, hence $b=3$. Thus $g=5$.

Now
\[
S^+=\Lambda^0E\oplus\Lambda^2E\oplus\Lambda^4E,
\]
and for $E=E_3\oplus E_2$,
\[
\Lambda^2E
=
\Lambda^2E_3
\oplus
(E_3\otimes E_2)
\oplus
\Lambda^2E_2,
\]
\[
\Lambda^4E
=
(\Lambda^3E_3\otimes E_2)
\oplus
(\Lambda^2E_3\otimes \Lambda^2E_2).
\]
These are exactly the six sectors
\[
(1,1)_0,\ (3,2)_{1/6},\ (\bar 3,1)_{-2/3},\ (\bar 3,1)_{1/3},\ (1,2)_{-1/2},\ (1,1)_1.
\]
Hence
\[
\dim S^+=2^{5-1}=16.
\]
Finally
\[
3q_-+2q_+=0
\]
with primitive integers and $q_-<0<q_+$ has the unique solution
\[
(q_-,q_+)=(-2,3).
\]
\end{proof}

\begin{corollary}[One chiral family emerges]
Under the hypotheses of \TFPTxref{thm:carrier-minimality-k}, the positive half-spinor packet $S^+$
carries exactly one chiral Standard Model family including $\nu^c$.
\end{corollary}

\begin{proof}
The six-sector decomposition displayed in \TFPTxref{thm:carrier-minimality-k} is exactly
\[
(1,1)_0\oplus(3,2)_{1/6}\oplus(\bar 3,1)_{-2/3}\oplus(\bar 3,1)_{1/3}\oplus(1,2)_{-1/2}\oplus(1,1)_1,
\]
namely one chiral Standard Model family together with $\nu^c$.
\end{proof}

\begin{remark}[Logical status of the carrier normal form]
The lemma above is the algebraic normal form of the former open carrier target Theorem~K. It is a rigidity
statement under two exterior signatures:
\[
\dim \Lambda^2 E_+=1,
\qquad
\Lambda^2 E_- \cong E_-^\vee.
\]
Within those signatures the packet size
\[
\dim\Lambda^{\mathrm{even}}E=16
\]
and the one-family decomposition are outputs, while the occupancy identities move behind the
carrier closure as derived consistency statements. On the retained branch, the compact Higgs index
first fixes the bosonic rank and \TFPTxref{thm:yukawa-forces-32,cor:yukawa-discharges-kpp} then
discharge these carrier-signature premises internally from the actual branch Yukawa tensor.
\end{remark}

\begin{corollary}[Carrier norm]
On the minimal carrier one has
\[
Y
=
\mathrm{diag}\!\left(-\frac13,-\frac13,-\frac13,\frac12,\frac12\right),
\qquad
\gamma=\tr_EY^2=\frac56.
\]
\end{corollary}

Hence from now on
\[
E=E_3\oplus E_2,
\qquad
X=-2P_3+3P_2,
\qquad
P_3=\frac{\mathbf 1-\epscar}{2},
\qquad
P_2=\frac{\mathbf 1+\epscar}{2},
\]
and the hypercharge operator is the normalized primitive two-point generator
\begin{equation}
\label{eq:explicit-hypercharge}
Y
=
\mathrm{diag}\!\left(-\frac13,-\frac13,-\frac13,\frac12,\frac12\right)
=
\frac{X}{6}
=
-\frac13 P_3+\frac12 P_2
=
\frac1{12}\mathbf 1+\frac5{12}\epscar.
\end{equation}

\begin{lemma}[Charge-generator / carrier-polarization inversion]
On the rigid carrier, the carrier polarization and the primitive two-point generator are recovered by
\begin{align}
\epscar&=\frac{2X-\mathbf 1}{5}=\frac{12Y-\mathbf 1}{5},
\label{eq:hypercharge-seam-inversion}\\
\qquad
X&=\frac{5\epscar+\mathbf 1}{2},
\nonumber\\
(2X-\mathbf 1)^2&=25\,\mathbf 1,
\nonumber\\
X^2-X-6\,\mathbf 1&=0,
\nonumber\\
6Y^2-Y-\mathbf 1&=0.
\label{eq:carrier-polynomial}
\end{align}
Consequently the carrier projectors are
\[
P_2=\frac{\mathbf 1+\epscar}{2}=\frac{2(3Y+\mathbf 1)}{5},
\qquad
P_3=\frac{\mathbf 1-\epscar}{2}=\frac{3(\mathbf 1-2Y)}{5}.
\]
\end{lemma}

\begin{proof}
The affine relations between $(X,Y)$ and $\epscar$ are inverted algebraically. Squaring the first
relation and using $\epscar^2=\mathbf 1$ gives the quadratic generator polynomial
$X^2-X-6\,\mathbf 1=0$, and dividing by $36$ yields the hypercharge polynomial. The projector
formulas then follow from $P_2=(\mathbf 1+\epscar)/2$ and $P_3=(\mathbf 1-\epscar)/2$.
\end{proof}

\begin{lemma}[Primitive trace-balanced carrier]
\label{lem:primitive-trace-balanced-carrier}
Let $Y$ be an endomorphism of a finite-dimensional carrier $E$ such that
\[
6Y^2-Y-\mathbf 1=0,
\qquad
\tr_EY=0,
\]
and assume the realization is primitive in the sense that $E$ is not a nontrivial direct sum of
identical trace-balanced $Y$-modules. Then
\[
\operatorname{Spec}(Y)=\left\{-\frac13,\frac12\right\},
\qquad
\dim\ker\!\left(Y+\frac13\mathbf 1\right)=3,
\qquad
\dim\ker\!\left(Y-\frac12\mathbf 1\right)=2.
\]
Equivalently,
\[
\rank E=5,
\qquad
E=E_3\oplus E_2,
\qquad
Y=-\frac13P_3+\frac12P_2.
\]
\end{lemma}

\begin{proof}
The quadratic polynomial factors as
\[
(2Y-\mathbf 1)(3Y+\mathbf 1)=0,
\]
so the only eigenvalues are
\[
y_+=\frac12,
\qquad
y_-=-\frac13.
\]
Write
\[
m_+:=\dim\ker(Y-y_+\mathbf 1),
\qquad
m_-:=\dim\ker(Y-y_-\mathbf 1).
\]
Then
\[
0=\tr_EY=\frac12\,m_+-\frac13\,m_-,
\]
so
\[
3m_+=2m_-.
\]
Because $\gcd(2,3)=1$, every positive integer solution has the form
\[
(m_+,m_-)=(2k,3k).
\]
Primitivity forces $k=1$, hence $(m_+,m_-)=(2,3)$. This is exactly the rigid $3+2$ split, with
the displayed spectral form of $Y$.
\end{proof}

\begin{remark}[Compressed carrier decoder]
The carrier can therefore be entered from the shorter algebraic normal form
\[
6Y^2-Y-\mathbf 1=0,
\qquad
\tr_EY=0,
\qquad
\text{primitive realization}
\Longrightarrow
E=E_3\oplus E_2.
\]
The exterior-signature route of \TFPTxref{thm:carrier-minimality-k} remains the theorem-level discharge
of these algebraic hypotheses from the actual branch geometry.
\end{remark}

The hard carrier package is summarized visually in \Cref{fig:minimal-carrier-glance}. The point is
that the same minimal data are readable in four equivalent ways: geometrically as an oriented
double cover with deck reflection and boundary polarization, algebraically as the trace-balanced
quadratic carrier decoder, combinatorially as the rigid $3+2$ split, and spectrally as the
primitive two-point generator $X$ with normalized hypercharge $Y=X/6$.

\begin{figure}[t]
\centering
\footnotesize
\begin{tikzpicture}[
    scale=0.88,
    transform shape,
    >=Latex,
    slot3/.style={draw=blue!60!black, fill=blue!6, rounded corners, minimum width=0.72cm, minimum height=0.72cm, align=center},
    slot2/.style={draw=orange!70!black, fill=orange!9, rounded corners, minimum width=0.72cm, minimum height=0.72cm, align=center},
    info/.style={draw=black!55, rounded corners, fill=gray!3, align=center, inner sep=5pt}
]
\node[font=\bfseries] at (-3.5,1.95) {oriented double cover};
\draw[draw=gray!70!black, fill=gray!6, rounded corners] (-5.4,0.55) rectangle (-1.7,1.35);
\draw[draw=gray!70!black, fill=gray!6, rounded corners] (-5.4,-1.35) rectangle (-1.7,-0.55);
\node at (-3.55,0.95) {upper sheet};
\node at (-3.55,-0.95) {lower sheet};
\draw[very thick, purple!70!black] (-5.6,0) -- (-1.45,0);
\node[anchor=west, purple!70!black] at (-1.38,0) {seam $\Sigma$};
\draw[<->, thick] (-3.55,0.42) -- (-3.55,-0.42);
\node[anchor=west] at (-3.4,0) {$\iota$};

\node[font=\bfseries] at (2.35,1.95) {rigid carrier split};
\node[slot3] (c1) at (0.55,0.95) {$c_1$};
\node[slot3, right=2mm of c1] (c2) {$c_2$};
\node[slot3, right=2mm of c2] (c3) {$c_3$};
\node[slot2, right=6mm of c3] (w1) {$w_1$};
\node[slot2, right=2mm of w1] (w2) {$w_2$};
\node[align=center, blue!60!black] at ($(c2.south)+(0,-0.45)$) {$E_3=P_3E$\\three color slots};
\node[align=center, orange!75!black] at ($(w1.south)+(0.45,-0.45)$) {$E_2=P_2E$\\two weak slots};
\draw[decorate, decoration={brace, amplitude=4pt}] ($(c1.north west)+(-0.07,0.12)$) -- ($(c3.north east)+(0.07,0.12)$);
\node[blue!60!black] at ($(c2.north)+(0,0.42)$) {$3$};
\draw[decorate, decoration={brace, amplitude=4pt}] ($(w1.north west)+(-0.07,0.12)$) -- ($(w2.north east)+(0.07,0.12)$);
\node[orange!75!black] at ($(w1.north east)+(0.2,0.42)$) {$2$};

\node[info] at (-0.45,-2.3) {%
    $\displaystyle
    X=-2P_3+3P_2,
    \qquad
    Y=\frac{X}{6}
    =-\frac13 P_3+\frac12 P_2
    =\frac1{12}\mathbf 1+\frac5{12}\iota
    \qquad
    \iota=\frac{2X-\mathbf 1}{5}
    $\\[2pt]
    $\displaystyle
    P_3=\frac{\mathbf 1-\iota}{2},
    \qquad
    P_2=\frac{\mathbf 1+\iota}{2},
    \qquad
    \mathrm{Spec}(Y)=\left\{-\frac13,\frac12\right\}
    $};
\end{tikzpicture}
\caption{Rigid carrier at a glance. The seam involution, the primitive two-point generator, and the
normalized hypercharge operator carry the same hard information: geometrically they separate the
two sheets, and algebraically they separate the $3+2$ carrier split through the projectors $P_3$
and $P_2$.}
\label{fig:minimal-carrier-glance}
\end{figure}

\begin{remark}[Two-point carrier algebra]
Because the minimal carrier satisfies $6Y^2-Y-\mathbf 1=0$, every analytic function of $Y$ reduces to
\[
f(Y)=A_f\,\mathbf 1+B_f\,Y,
\qquad
A_f=\frac{2f(1/2)+3f(-1/3)}{5},
\qquad
B_f=\frac{6\bigl(f(1/2)-f(-1/3)\bigr)}{5}.
\]
Thus the hard carrier is effectively a two-point algebra supported on the spectral values $-1/3$ and $1/2$, and the projectors $P_3,P_2$ are already polynomial functions of $Y$.
\end{remark}

\begin{remark}[Factorized carrier polynomial and carrier discriminant]\label{rem:carrier-discriminant}
For a general split carrier $E=E_b\oplus E_s$ with total rank $g:=b+s$ and eigenvalues $-1/b$ and $1/s$, the carrier generator satisfies
\[
\left(Y+\frac{1}{b}\right)\left(Y-\frac{1}{s}\right)=0
\quad\Longleftrightarrow\quad
bs\,Y^2+(s-b)Y-\mathbf 1=0.
\]
Its discriminant is
\[
\Delta_Y=(s-b)^2+4bs=(b+s)^2=g^2.
\]
Thus the square root of the carrier discriminant is the carrier rank itself. Moreover,
\[
Y=\frac{b-s}{2bs}\,\mathbf 1+\frac{g}{2bs}\,\epscar,
\qquad
\epscar=\frac{2bs\,Y-(b-s)\,\mathbf 1}{g}.
\]
On the minimal branch $(b,s)=(3,2)$, this reduces to
\[
6Y^2-Y-\mathbf 1=(2Y-\mathbf 1)(3Y+\mathbf 1),
\qquad
\Delta_Y=25=5^2,
\qquad
\epscar=\frac{12Y-\mathbf 1}{5}.
\]
Hence the denominator $5$ in the hypercharge--seam inversion is exactly $\sqrt{\Delta_Y}$, not an additional inserted constant.

This does not replace the stabilizer proof of
\[
G_{\mathrm{car}}=S(U(3)\times U(2)),
\]
which is the actual source of the internal carrier stabilizer theorem in the present manuscript.
What the factorization and discriminant do show exactly is that the hard carrier already supports the distinguished spectral values
\[
Y=\frac{1}{2},
\qquad
Y=-\frac{1}{3},
\]
and that the inversion denominator is fixed algebraically by the carrier discriminant.
\end{remark}

The two carrier invariants used repeatedly below are
\[
\gamma=\tr_EY^2=\frac56,
\qquad
B=\frac{\rank E_3}{\rank E_2}=\frac32.
\]
\begin{corollary}[Carrier norm from the carrier polynomial]
The carrier polynomial already fixes
\[
\gamma=\tr_EY^2=\frac16\tr_E(Y+\mathbf 1)=\frac56,
\]
since $\tr_EY=0$ and $\rank E=5$.
\end{corollary}

\begin{corollary}[Minimal carrier compression]
\label{cor:minimal-carrier-compression}
On the minimal carrier branch,
\[
B\gamma=\frac32\cdot\frac56=\frac54.
\]
\end{corollary}

We package the hard kernel as
\[
\Cmin=(E_3\oplus E_2,\ X,\ Y,\ S^+,\ \gamma,\ B).
\]

\subsection{Internal carrier stabilizer and one-family packet}

\begin{theorem}[Internal carrier stabilizer]
\label{thm:global-sm-theorem}
The determinant-preserving connected stabilizer of the minimal carrier split is
\[
G_{\mathrm{car}}
=
S(U(3)\times U(2))
\cong
\frac{SU(3)\times SU(2)\times U(1)_Y}{\mathbb Z_6}.
\]
The corresponding Lie algebra factorization
\[
\mathfrak g_{\SM}=\mathfrak{su}(3)\oplus\mathfrak{su}(2)\oplus\mathfrak u(1)
\]
is only the differential, local form of the same statement; throughout the manuscript this quotient
is intended as the internal carrier-stabilizer form whenever a group symbol is written, and the
local algebra notation is used only when explicitly stated.
\end{theorem}

\begin{proof}
An automorphism preserving the split $E_3\oplus E_2$ acts by a pair
$(u_3,u_2)\in U(3)\times U(2)$. Requiring determinant preservation on the full
carrier gives
\[
\det(u_3)\det(u_2)=1,
\]
so the connected stabilizer is exactly $S(U(3)\times U(2))$.

Define
\[
\Phi:SU(3)\times SU(2)\times U(1)_Y \longrightarrow S(U(3)\times U(2)),
\qquad
\Phi(g_3,g_2,z)=\bigl(z^2 g_3,\ z^{-3} g_2\bigr).
\]
This is a well-defined group homomorphism because the exponents are integers, and
\[
\det(z^2 g_3)\det(z^{-3} g_2)=z^6 z^{-6}=1.
\]

To prove surjectivity, take $(u_3,u_2)\in S(U(3)\times U(2))$. Choose
$z\in U(1)$ with
\[
z^6=\det(u_3)=\det(u_2)^{-1}.
\]
Then
\[
g_3:=z^{-2}u_3\in SU(3),
\qquad
g_2:=z^{3}u_2\in SU(2),
\]
and therefore
\[
\Phi(g_3,g_2,z)=(u_3,u_2).
\]

If $\Phi(g_3,g_2,z)=(\mathbf 1_3,\mathbf 1_2)$, then
\[
g_3=z^{-2}\mathbf 1_3,
\qquad
g_2=z^{3}\mathbf 1_2.
\]
The conditions $g_3\in SU(3)$ and $g_2\in SU(2)$ imply $z^6=1$. Hence
\[
\ker\Phi
=
\left\{
\bigl(z^{-2}\mathbf 1_3,\ z^{3}\mathbf 1_2,\ z\bigr): z^6=1
\right\}
\cong \mathbb Z_6.
\]
Therefore
\[
S(U(3)\times U(2))
\cong
\frac{SU(3)\times SU(2)\times U(1)_Y}{\mathbb Z_6}.
\]
\end{proof}

\begin{theorem}[Faithful physical global gauge group]
\label{thm:physical-gauge-group}
\label{thm:physical-gauge-group-faithful}
The action of
\[
S(U(3)\times U(2))
\]
on the admissible matter-plus-Higgs package
\[
\mathcal R_{\mathrm{adm}}:=S^+\oplus E_2
\]
is faithful. Therefore
\[
G_{\mathrm{phys}}=G_{\mathrm{car}}=S(U(3)\times U(2)).
\]
Equivalently,
\[
G_{\mathrm{phys}}
\cong
\frac{SU(3)\times SU(2)\times U(1)_Y}{\mathbb Z_6}.
\]
No further nontrivial quotient survives on the physical matter-plus-Higgs package.
\end{theorem}

\begin{proof}
Let $(u_3,u_2)\in S(U(3)\times U(2))$ act trivially on $\mathcal R_{\mathrm{adm}}$. The Higgs
sector is $E_2$, so triviality there implies
\[
u_2=\mathbf 1_2.
\]
The matter packet contains the quark doublet sector
\[
Q=E_3\otimes E_2.
\]
Since $u_2=\mathbf 1_2$, trivial action on $Q$ forces
\[
u_3=\mathbf 1_3.
\]
Hence the kernel is trivial. The quotient from $SU(3)\times SU(2)\times U(1)$ to
$S(U(3)\times U(2))$ already has kernel $\mathbb Z_6$ by \TFPTxref{thm:global-sm-theorem}, so this is
the exact physical global form.
\end{proof}

Throughout the manuscript, $S(U(3)\times U(2))$ is therefore read not only as the internal carrier
stabilizer but as the physical global gauge group selected by the faithful action on the admissible
matter-plus-Higgs package.

Matter is read from the positive half-spinor
\[
S^+=\Lambda^{\mathrm{even}}E,
\qquad
\dim S^+=16.
\]

\begin{theorem}[One carrier family from the exterior packet]
The positive half-spinor of the minimal carrier decomposes as
\[
S^+
\cong
(1,1)_0\oplus(3,2)_{1/6}\oplus(\bar 3,1)_{-2/3}\oplus(\bar 3,1)_{1/3}\oplus(1,2)_{-1/2}\oplus(1,1)_1,
\]
that is, exactly one chiral Standard Model family with $\nu^c$. Equivalently, in the familiar unified packaging language,
\[
S^+\cong \mathbf 1\oplus \mathbf{10}\oplus \overline{\mathbf 5}.
\]
\end{theorem}

\begin{proof}
One has
\[
S^+=\Lambda^0E\oplus\Lambda^2E\oplus\Lambda^4E.
\]
The $\Lambda^0E$ sector is immediately
\[
\Lambda^0E\cong (1,1)_0,
\]
which is identified with $\nu^c$ in left-chiral notation. For the minimal carrier,
\[
\Lambda^2E
\cong
\Lambda^2E_3\oplus(E_3\otimes E_2)\oplus\Lambda^2E_2.
\]
Here
\[
\Lambda^2E_3\cong \bar 3,
\qquad
E_3\otimes E_2\cong (3,2),
\qquad
\Lambda^2E_2\cong 1,
\]
and the hypercharges are obtained by adding the carrier charges from
\[
Y=\mathrm{diag}\!\left(-\frac13,-\frac13,-\frac13,\frac12,\frac12\right).
\]
Therefore
\[
\Lambda^2E
\cong
(\bar 3,1)_{-2/3}\oplus(3,2)_{1/6}\oplus(1,1)_1.
\]
Because the full determinant is trivial in $S(U(3)\times U(2))$, one has $\Lambda^5E\cong 1$ and hence
\[
\Lambda^4E\cong E^\ast.
\]
Equivalently,
\[
\Lambda^4E
\cong
(\Lambda^3E_3\otimes E_2)\oplus(\Lambda^2E_3\otimes\Lambda^2E_2).
\]
Using $\Lambda^3E_3\cong 1$, $\Lambda^2E_2\cong 1$, and $\Lambda^2E_3\cong E_3^\ast\cong\bar 3$, one gets
\[
\Lambda^4E\cong (1,2)_{-1/2}\oplus(\bar 3,1)_{1/3}.
\]
Collecting the three even-degree sectors yields
\[
S^+
\cong
(1,1)_0\oplus(\bar 3,1)_{-2/3}\oplus(3,2)_{1/6}\oplus(1,1)_1\oplus(\bar 3,1)_{1/3}\oplus(1,2)_{-1/2},
\]
with total multiplicity $1+3+6+1+3+2=16$. The $\Lambda^2E$ sector is the standard left-chiral $\mathbf{10}$ of $SU(5)$, the $\Lambda^4E$ sector is the standard left-chiral $\overline{\mathbf 5}$, and the $\Lambda^0E$ sector is the singlet $\mathbf 1$.
\end{proof}

In left-chiral notation the one-family decomposition takes the explicit form shown in \TFPTxref{tab:one-family}. This is the place where the manuscript must be explicit: the carrier theorem is meaningful only if the reader can see how the sixteen states are organized.

\begin{table}[t]
\centering
\footnotesize
\begin{tabular}{@{}>{\raggedright\arraybackslash}p{0.18\linewidth}>{\raggedright\arraybackslash}p{0.19\linewidth}>{\raggedleft\arraybackslash}p{0.13\linewidth}>{\raggedright\arraybackslash}p{0.20\linewidth}>{\raggedleft\arraybackslash}p{0.08\linewidth}@{}}
\toprule
Exterior component & $\SM$ representation & Hypercharge & Field content & Mult. \\
\midrule
$\Lambda^0E$ & $(1,1)_0$ & $0$ & $\nu^c$ & $1$ \\
$\Lambda^2E_3$ & $(\bar 3,1)_{-2/3}$ & $-2/3$ & $u^c$ & $3$ \\
$E_3\otimes E_2$ & $(3,2)_{1/6}$ & $1/6$ & $Q_L$ & $6$ \\
$\Lambda^2E_2$ & $(1,1)_1$ & $1$ & $e^c$ & $1$ \\
\makecell[l]{$\Lambda^3E_3$\\$\otimes E_2$} & $(1,2)_{-1/2}$ & $-1/2$ & $L_L$ & $2$ \\
\makecell[l]{$\Lambda^2E_3$\\$\otimes\Lambda^2E_2$} & $(\bar 3,1)_{1/3}$ & $1/3$ & $d^c$ & $3$ \\
\midrule
\multicolumn{4}{r}{Total multiplicity} & $16$ \\
\bottomrule
\end{tabular}
\caption{One chiral family from $S^+=\Lambda^{\mathrm{even}}E$, written in left-chiral notation with charge-conjugate singlets.}
\label{tab:one-family}
\end{table}

\begin{proposition}[Exact parity-code realization of the family packet]
\label{prop:parity-code-family}
Identifying each wedge monomial in $\Lambda E$ with its occupation word
\[
x=(x_1,\dots,x_5)\in\{0,1\}^5,
\]
the positive half-spinor is exactly the even-parity subspace
\[
S^+
\cong
\mathrm{span}\Bigl\{|x\rangle:\sum_{i=1}^5 x_i\equiv 0 \pmod 2\Bigr\}.
\]
Equivalently,
\[
S^+
=
\bigl(\Lambda^{\mathrm{even}}E_3\otimes\Lambda^{\mathrm{even}}E_2\bigr)
\oplus
\bigl(\Lambda^{\mathrm{odd}}E_3\otimes\Lambda^{\mathrm{odd}}E_2\bigr),
\]
and each summand has dimension $8$.
\end{proposition}

\begin{proof}
Choose a basis of $E$ adapted to the split $E_3\oplus E_2$. Exterior monomials are then in bijection with occupation words $x\in\{0,1\}^5$, and the exterior degree is exactly the Hamming weight $\sum_i x_i$. Hence $\Lambda^{\mathrm{even}}E$ is the span of the even-parity words. Splitting the degree into the $E_3$ and $E_2$ contributions gives the direct-sum decomposition above. Since
\[
\dim \Lambda^{\mathrm{even}}E_3=1+\binom32=4,
\qquad
\dim \Lambda^{\mathrm{odd}}E_3=\binom31+\binom33=4,
\]
and
\[
\dim \Lambda^{\mathrm{even}}E_2=1+\binom22=2,
\qquad
\dim \Lambda^{\mathrm{odd}}E_2=\binom21=2,
\]
each tensor-product sector has dimension $8$.
\end{proof}

\begin{remark}[Exact parity lock, not a loose code analogy]
The family packet is therefore not merely code-like: at the level of occupation kinematics it is exactly the classical even-parity $[5,4,2]$ code on five bits. What the present paper still does \emph{not} claim is a dynamical error-correction theorem for realistic noise channels. The exact statement is the parity lock itself. In the maximally mixed state on the $16$ occupation words, eight states lie in the $(\mathrm{even},\mathrm{even})$ sector and eight in $(\mathrm{odd},\mathrm{odd})$, so the coarse parity observables of the $3$-side and $2$-side carry exactly one bit of mutual information.
\end{remark}

The family packet can also be read as a two-dimensional occupancy map, shown in \Cref{fig:family-occupancy-map}. This is where the one-family table, the even-parity code, the split weight enumerator, and the $X=6Y$ arithmetic all meet in one glance.

\begin{figure}[t]
\centering
\footnotesize
\begin{tikzpicture}[>=Latex, x=1.7cm, y=1.45cm]
\draw[->, thick] (-0.2,0) -- (3.4,0) node[below] {$n_3$};
\draw[->, thick] (0,-0.2) -- (0,2.35) node[left] {$n_2$};
\foreach \x in {0,1,2,3} {
    \draw[gray!30] (\x,0) -- (\x,2.05);
    \node[below] at (\x,-0.02) {\x};
}
\foreach \y in {0,1,2} {
    \draw[gray!30] (0,\y) -- (3.05,\y);
    \node[left] at (-0.02,\y) {\y};
}
\node[align=left, anchor=west] at (2.05,2.15) {$n_3+n_2$ even\\$\Longleftrightarrow S^+$ admissible};

\foreach \x/\y in {1/0,3/0,0/1,2/1,1/2,3/2} {
    \draw[red!65!black, line width=0.6pt] (\x-0.08,\y-0.08) -- (\x+0.08,\y+0.08);
    \draw[red!65!black, line width=0.6pt] (\x-0.08,\y+0.08) -- (\x+0.08,\y-0.08);
}

\filldraw[blue!55!black] (0,0) circle (3.8pt);
\node[text=white, font=\scriptsize] at (0,0) {1};
\node[anchor=west] at (0.14,0.05) {$\nu^c$};

\filldraw[blue!55!black] (2,0) circle (4.5pt);
\node[text=white, font=\scriptsize] at (2,0) {3};
\node[anchor=west] at (2.12,0.08) {$u^c$};

\filldraw[green!55!black] (1,1) circle (5.5pt);
\node[text=white, font=\scriptsize] at (1,1) {6};
\node[anchor=west] at (1.14,1.08) {$Q_L$};

\filldraw[orange!80!black] (0,2) circle (3.8pt);
\node[text=white, font=\scriptsize] at (0,2) {1};
\node[anchor=west] at (0.12,1.9) {$e^c$};

\filldraw[orange!80!black] (3,1) circle (4.3pt);
\node[text=white, font=\scriptsize] at (3,1) {2};
\node[anchor=west] at (3.12,1.03) {$L_L$};

\filldraw[orange!80!black] (2,2) circle (4.5pt);
\node[text=white, font=\scriptsize] at (2,2) {3};
\node[anchor=west] at (2.12,1.9) {$d^c$};

\node[draw=black!55, rounded corners, fill=gray!4, align=center] at (1.55,-0.72) {%
    multiplicity $=\binom{3}{n_3}\binom{2}{n_2}$\\
    scaled hypercharge $X=3n_2-2n_3=6Y$};
\end{tikzpicture}
\caption{Occupancy map of the family packet. The admissible sectors are exactly the even lattice points in $(n_3,n_2)$, with multiplicities $(1,3,6,1,2,3)$ and field labels matching \TFPTxref{tab:one-family}. This is the visual meeting point of parity lock, split weight enumerator, and family hypercharge arithmetic.}
\label{fig:family-occupancy-map}
\end{figure}

\subsection{Character, trace identities, and the abelian index}

\begin{proposition}[Split weight enumerator and one-family character]
\label{prop:split-weight-enumerator}
For the present paragraph let
\[
\chi_{S^+}(t):=\tr_{S^+}(e^{tY}).
\]
Define also the split weight enumerator
\[
W(u,v)
:=
\sum_{\substack{0\le n_3\le 3\\ 0\le n_2\le 2\\ n_3+n_2\ \mathrm{even}}}
\binom{3}{n_3}\binom{2}{n_2}u^{n_3}v^{n_2}.
\]
Then
\begin{equation}
\label{eq:split-weight-enumerator}
W(u,v)
=
1+3u^2+6uv+v^2+2u^3v+3u^2v^2
=
\frac12\Big[(1+u)^3(1+v)^2+(1-u)^3(1-v)^2\Big].
\end{equation}
Moreover,
\begin{equation}
\label{eq:family-character}
\chi_{S^+}(t)
=
\frac12\Big[
\bigl(1+e^{-t/3}\bigr)^3\bigl(1+e^{t/2}\bigr)^2
+
\bigl(1-e^{-t/3}\bigr)^3\bigl(1-e^{t/2}\bigr)^2
\Big]
=
W(e^{-t/3},e^{t/2}).
\end{equation}
Consequently,
\[
\chi_{S^+}(0)=16,
\qquad
\chi_{S^+}'(0)=\tr_{S^+}Y=0,
\qquad
\chi_{S^+}''(0)=\tr_{S^+}Y^2=\frac{10}{3},
\qquad
\chi_{S^+}'''(0)=\tr_{S^+}Y^3=0.
\]
\end{proposition}

\begin{remark}[Master generating function for the one-family packet]
The even-parity split weight enumerator
\[
W(u,v)=\frac12\bigl[(1+u)^3(1+v)^2+(1-u)^3(1-v)^2\bigr]
\]
is the master generator of the family packet. The one-family character is obtained by the
specialization $u=e^{-t/3}$ and $v=e^{t/2}$, while all trace moments of $Y$ are recovered by
differentiating at $t=0$.
\end{remark}

\begin{proof}
The pairs $(n_3,n_2)$ allowed by even total degree are
\[
(0,0),\ (2,0),\ (1,1),\ (0,2),\ (3,1),\ (2,2),
\]
with multiplicities
\[
1,\ 3,\ 6,\ 1,\ 2,\ 3.
\]
This gives the explicit polynomial formula for $W(u,v)$. Equivalently, the parity projector
\[
\frac12\bigl(1+(-1)^{n_3+n_2}\bigr)
\]
acting on $(1+u)^3(1+v)^2$ yields \Cref{eq:split-weight-enumerator}. Substituting
\[
u=e^{-t/3},
\qquad
v=e^{t/2}
\]
gives \Cref{eq:family-character}. For an even exterior algebra one also has
\[
\tr_{\Lambda^{\mathrm{even}}E}(e^{tY})
=
\frac12\Big[\prod_{j=1}^5(1+e^{tq_j})+\prod_{j=1}^5(1-e^{tq_j})\Big],
\]
with carrier charges $q_j\in\{-1/3,-1/3,-1/3,1/2,1/2\}$. Differentiating at $t=0$ yields the stated dimension and trace identities.
\end{proof}

\begin{corollary}[Multiplicative character as six-sector decoder]
\label{cor:multiplicative-family-character}
Let
\[
X:=6Y.
\]
Then the multiplicative one-family character
\[
\Xi_{S^+}(t):=\tr_{S^+}(t^X)
\]
is
\begin{equation}
\label{eq:multiplicative-family-character}
\Xi_{S^+}(t)
=
\frac12\Big[
(1+t^{-2})^3(1+t^3)^2
+
(1-t^{-2})^3(1-t^3)^2
\Big]
=
t^6+3t^2+6t+1+2t^{-3}+3t^{-4}.
\end{equation}
Consequently
\[
\Xi_{S^+}(1)=16,
\qquad
(t\partial_t)\Xi_{S^+}(1)=\tr_{S^+}X=0,
\qquad
(t\partial_t)^2\Xi_{S^+}(1)=\tr_{S^+}X^2=120=5!.
\]
Thus the six sectors and their multiplicities are read off directly from the monomials of
\Cref{eq:multiplicative-family-character}.
\end{corollary}

\begin{proof}
Since $X=3N_2-2N_3$, every admissible occupancy pair contributes the monomial
$t^{3n_2-2n_3}$. Therefore
\[
\Xi_{S^+}(t)=W(t^{-2},t^3),
\]
and \Cref{eq:split-weight-enumerator} gives the displayed factorized form. Expanding yields the six
monomials shown above. Applying $t\partial_t$ once and twice at $t=1$ recovers the first two trace
moments of $X$.
\end{proof}

\begin{remark}[Table as readout, character as proof device]
The one-family table remains useful as a reader-facing decoding sheet, but the whole six-sector
arithmetic is already compressed into the single multiplicative character
\Cref{eq:multiplicative-family-character}.
\end{remark}

\begin{remark}[Six admissible occupancy sectors]
The six monomials in \Cref{eq:split-weight-enumerator} are exactly the six allowed occupancy pairs $(n_3,n_2)$ of the family packet. This is the shortest combinatorial route from the split carrier to the six Standard Model multiplet sectors listed in \TFPTxref{tab:one-family}.
\end{remark}

The carrier packet therefore reproduces the standard trace cancellations
\[
\tr_{S^+}Y=0,
\qquad
\tr_{S^+}Y^3=0,
\]
and packages them together with the quadratic trace in one formula. Since $\tr_{S^+}Y^2=4\gamma=10/3$ and $\tr_\Phi Y^2=1/2$ for one weak doublet with hypercharge $1/2$, the hypercharge trace coefficient organizes as
\begin{equation}
\label{eq:abelian-index}
\bone(N_{\mathrm{fam}},N_\Phi)
=
\frac25N_{\mathrm{fam}}\tr_{S^+}Y^2+\frac15N_\Phi\tr_\Phi Y^2
=
\frac85N_{\mathrm{fam}}\gamma+\frac1{10}N_\Phi
=
\frac43N_{\mathrm{fam}}+\frac1{10}N_\Phi.
\end{equation}
On the canonical branch $(N_{\mathrm{fam}},N_\Phi)=(3,1)$ this gives
\[
\bone=\frac{41}{10}.
\]
The abelian index is therefore not inserted as an isolated Standard Model datum; it is a carrier-level output once the family and Higgs counts are declared.

\begin{remark}[Canonical-branch Diophantine lock]
If one imposes $\bone=41/10$ on nonnegative integer data, then
\[
40N_{\mathrm{fam}}+3N_\Phi=123.
\]
The physically nontrivial solution is $(N_{\mathrm{fam}},N_\Phi)=(3,1)$, so the canonical branch is already sharply constrained at the level of the trace arithmetic.
\end{remark}

\begin{proposition}[Family hypercharge arithmetic and compression]
\label{prop:family-hypercharge-compression}
Let $N_3$ and $N_2$ be the occupation-number operators on the $E_3$ and $E_2$ slots, restricted to $S^+$, and define the scaled hypercharge operator
\[
X:=6Y=3N_2-2N_3.
\]
Then on $S^+$ one has
\begin{align}
X&\equiv N_2\equiv N_3 \pmod 2,\label{eq:x-mod-2}\\
X&\equiv N_3 \pmod 3,\label{eq:x-mod-3}\\
X&\equiv 3(N_2+N_3) \pmod 5.\label{eq:x-mod-5}
\end{align}
In particular, the even exterior degree $k=N_2+N_3\in\{0,2,4\}$ is read off from the residue class of $X$ modulo $5$: if $X\equiv r\pmod 5$ with $r\in\{0,1,2\}$, then $k=2r$, so
\[
X\bmod 5=0,1,2
\qquad\Longleftrightarrow\qquad
\Lambda^0E,\ \Lambda^2E,\ \Lambda^4E.
\]
Moreover,
\[
\mathrm{Spec}(X|_{S^+})=\{0,-4,1,6,-3,2\},
\]
and therefore
\begin{equation}
\label{eq:family-x-polynomial}
X(X+4)(X-1)(X-6)(X+3)(X-2)=0
\qquad\text{on }S^+.
\end{equation}
\end{proposition}

\begin{proof}
Each occupied $E_3$ mode contributes hypercharge $-1/3$ and each occupied $E_2$ mode contributes $+1/2$, so multiplying by $6$ gives $X=3N_2-2N_3$. Modulo $2$, the term $-2N_3$ vanishes, so $X\equiv N_2$. Because $S^+$ has even total degree, $N_2\equiv N_3\pmod 2$, which proves \Cref{eq:x-mod-2}. Modulo $3$, one has $3N_2\equiv 0$ and $-2\equiv 1$, giving \Cref{eq:x-mod-3}. Modulo $5$, the congruence $-2\equiv 3$ yields \Cref{eq:x-mod-5}. Since on $S^+$ the even degree is $k=N_2+N_3\in\{0,2,4\}$, the possible residues are $3k\equiv 0,1,2\pmod 5$, from which $k=2r$ with $r=X\bmod 5$ follows. Evaluating $X=3n_2-2n_3$ on the six admissible occupancy pairs $(n_3,n_2)$ gives the listed spectrum. The polynomial \Cref{eq:family-x-polynomial} is then the spectral polynomial with these six simple roots.
\end{proof}

\begin{corollary}[Five-fold phase filters for even degree]
\label{cor:family-degree-filters}
Let
\[
\zeta_5=e^{2\pi i/5},
\qquad
U_5:=e^{2\pi iX/5}.
\]
Then the even-degree projectors are the discrete Fourier filters
\[
P_{\Lambda^{2m}E}
=
\frac15\sum_{r=0}^{4}\zeta_5^{-mr}U_5^r,
\qquad
m=0,1,2.
\]
\end{corollary}

\begin{proof}
By \TFPTxref{prop:family-hypercharge-compression}, the eigenspaces $\Lambda^{2m}E$ carry the residue class $X\equiv m\pmod 5$. Hence $U_5$ has eigenvalue $\zeta_5^m$ on $\Lambda^{2m}E$, and the standard Fourier projector formula on the cyclic group $\mathbb Z_5$ gives the claim.
\end{proof}

\begin{remark}[From the two-point carrier to the six-point family algebra]
The minimal carrier itself is a two-point algebra, but the family packet lifts the scaled hypercharge to a six-point algebra. Therefore every analytic function of $X$ on $S^+$ reduces to a polynomial of degree at most five. If
\[
M_n:=\tr_{S^+}(X^n),
\]
then \Cref{eq:family-x-polynomial} implies the finite recursion
\[
M_{n+6}
=
2M_{n+5}+31M_{n+4}-20M_{n+3}-156M_{n+2}+144M_{n+1}.
\]
In this sense the whole hypercharge statistics of one family is finitely compressed into six spectral values.
\end{remark}


% ---- Compression identities and carrier-relevant comparison setup: source lines 2696-3156 ----
\section{Compression identities and numerical comparison of alternative discrete worlds}

This section collects the algebraic compressions that turn the carrier packet and kernel numerics
from a list of internally consistent outputs into an overdetermined structure. These compression
identities are derived consistency statements on the retained branch: their formal derivation closes
only once the bosonic-rank and branch-Yukawa carrier theorems have fixed the rigid $3+2$ split. By
contrast, the discrete-world table below is a numerical comparison layer: it is not part of the
rigid theorem surface and should be read as an appendix-style pressure test on the carrier packet.

\subsection{Derived \texorpdfstring{$48=4\cdot12$}{48=4·12} identity and the \texorpdfstring{$2$-$3$-$5$}{2-3-5} arithmetic}

\begin{proposition}[Derived \texorpdfstring{$48=4\cdot12$}{48=4·12} identity on the retained branch]
\label{prop:rank-lock}
On the retained branch, once the family multiplicity and carrier ranks have been fixed, one
recovers
\[
\Oadm
=
N_{\mathrm{fam}}\,\dim S^+
=
(g-1)\dim\mathfrak g_{\mathrm{car}}
=
(g-1)(b^2+s^2-1)
\]
with
\[
N_{\mathrm{fam}}=3,
\qquad
\dim S^+=16,
\qquad
(b,s)=(3,2),
\qquad
g=5.
\]
Equivalently, the former rank-lock formula is a derived consistency identity evaluated at the
closed branch point:
\begin{equation}
\label{eq:rank-lock}
3\cdot 2^{g-1}=(g-1)(b^2+s^2-1).
\end{equation}
\end{proposition}

\begin{proof}
\Cref{thm:harmonic-family-modes} gives the physical family multiplicity
\[
N_{\mathrm{fam}}=\dim\mathcal H_F=3.
\]
The carrier closure is completed later by \TFPTxref{cor:bosonic-rank-two,thm:yukawa-forces-32}, which
fix
\[
(b,s)=(3,2),
\qquad
g=5,
\qquad
\dim S^+=16.
\]
Therefore
\[
\Oadm=N_{\mathrm{fam}}\dim S^+=3\cdot16=48.
\]
For the same minimal branch,
\[
\dim\mathfrak g_{\mathrm{car}}
=
\dim\mathfrak{su}(3)+\dim\mathfrak{su}(2)+\dim\mathfrak u(1)
=
12.
\]
Therefore
\[
\Oadm=48=(g-1)\dim\mathfrak g_{\mathrm{car}}=4\cdot 12,
\]
which is exactly the displayed derived identity written in branch variables. The former occupancy
balance is thus a consistency check rather than an input assumption.
\end{proof}

\begin{remark}[The $2$-$3$-$5$ arithmetic]
The hard carrier packet rests on exactly three discrete numbers: $2$ (the double cover), $3$ (the color rank and family count), and $5$ (the carrier rank). From these three atoms one recovers
\[
\gamma=\frac{5}{6}=\frac{5}{2\cdot 3},
\qquad
B=\frac{3}{2},
\qquad
B\gamma=\frac{5}{4}=\frac{5}{2^2},
\]
\[
\Oadm=3\cdot 2^4=48,
\qquad
\Dstart=5\cdot 12=60,
\qquad
\frac{\Dstart}{\Oadm}=\frac{5}{4}=B\gamma,
\]
and the defect coefficients $\deltatop=48\,\cthree^4$ and $\delta_2=\frac54\,\deltatop^2=2880\,\cthree^8$ with $2880=60\cdot 48$.
The persistent appearance of $5/4$ as carrier asymmetry $\times$ norm, second defect stage, and ladder-to-occupancy ratio is therefore not three coincidences but one rational shadow of the $2$-$3$-$5$ triad.
\end{remark}

\begin{lemma}[Universal \texorpdfstring{$5/4$}{5/4} compression quotient on the minimal branch]
\label{lem:universal-5-4-compression}
Define
\[
\rho_{5/4}:=B\gamma.
\]
On the canonical minimal branch,
\[
\rho_{5/4}
=
B\gamma
=
\frac{\delta_2}{\deltatop^2}
=
\frac{\Dstart}{\Oadm}
=
\frac{\gcar}{\gcar-1}
=
\frac54.
\]
Thus $5/4$ is the universal compression quotient of the minimal branch rather than a recurring numerical motif.
\end{lemma}

\begin{proof}
The identity $B\gamma=\frac54$ is \Cref{cor:minimal-carrier-compression}. The defect relation gives
\[
\frac{\delta_2}{\deltatop^2}=\frac54.
\]
On the canonical branch,
\[
\frac{\Dstart}{\Oadm}=\frac{60}{48}=\frac54,
\]
and with $\gcar=5$ one also has
\[
\frac{\gcar}{\gcar-1}=\frac{5}{4}.
\]
Hence all displayed ratios agree on the same exact value.
\end{proof}

\begin{remark}[Interpretive status of the quotient]
The quotient itself is exact. Any holographic or spacetime-dimensional reading of the denominator $\gcar-1=4$ remains interpretive in the present manuscript rather than theorem-level.
\end{remark}

\begin{corollary}[Rank-5 master compression on the retained branch]
\label{cor:rank-5-master-compression}
On the retained branch one has $\gcar=5$. Consequently
\[
N_{\mathrm{fam}}=\frac{\gcar+1}{2}=3,
\qquad
\Oadm=(\gcar+1)\,2^{\gcar-2}=48,
\]
\[
\Oadm\gamma=\gcar\,2^{\gcar-2}=40,
\qquad
\bone=\frac{\gcar\,2^{\gcar-2}+1}{10}=\frac{41}{10},
\qquad
10\,\bone=41.
\]
Thus the closed family count, admissible occupancy, weighted occupancy, and abelian index all
compress to the single carrier rank.
\end{corollary}

\begin{proof}
By \TFPTxref{cor:bosonic-rank-two,thm:yukawa-forces-32}, the retained carrier has
\[
\gcar=5.
\]
By \Cref{thm:harmonic-family-modes}, the family block fixes $N_{\mathrm{fam}}=3$, hence
\[
\Oadm=N_{\mathrm{fam}}\,2^{\gcar-1}=3\cdot 2^4=48.
\]
The carrier norm corollary gives $\gamma=5/6$, so
\[
\Oadm\gamma=48\cdot\frac56=40.
\]
Finally the closed Higgs branch has $N_\Phi=1$, and \TFPTeqref{eq:41-chain-main} yields
\[
10\,\bone=\Oadm\gamma+N_\Phi=40+1=41.
\]
\end{proof}

\begin{corollary}[Pythagorean compression of the abelian index on the retained branch]
\label{cor:pythagorean-integrity-index}
On the retained branch,
\[
\mathcal I_{41}
=
10\,\bone
=
2\gcar(\gcar-1)+1
=
\gcar^2+(\gcar-1)^2.
\]
In particular, since $\gcar=5$ on the retained branch, one gets
\[
\mathcal I_{41}=5^2+4^2=41.
\]
\end{corollary}

\begin{proof}
By \TFPTxref{cor:rank-5-master-compression}, one has $\gcar=5$ and
\[
10\,\bone=\gcar\,2^{\gcar-2}+1=5\cdot 2^3+1=41.
\]
\[
2\gcar(\gcar-1)+1=2\cdot 5\cdot 4+1=41.
\]
Since on the retained branch
\[
2\gcar(\gcar-1)+1=\gcar^2+(\gcar-1)^2,
\]
the claim follows.
\end{proof}

\begin{remark}[The number $4$ as a cross-cutting echo]
In addition to the $2$-$3$-$5$ triad, the integer $4$ appears at four structurally distinct levels:
\[
F_{\mathrm{patch}}=4,
\qquad
\gcar-1=4,
\qquad
\frac{\Oadm}{\dim\mathfrak g_{\SM}}=\frac{48}{12}=4,
\qquad
\frac{|R(E_8)|}{\Dstart}=\frac{240}{60}=4.
\]
That is, four corner patches, carrier rank minus one, matter-to-gauge compression factor, and root-to-start quotient all evaluate to the same integer.
\end{remark}

\subsection{Carrier compressions on the minimal branch}

\begin{corollary}[Two-point spectral compression of $\gamma$ and $\phibase$]
Let
\[
y_+:=\frac12,
\qquad
y_-:=-\frac13
\]
be the two distinct carrier eigenvalues singled out by the factorized carrier polynomial. Then
\[
\gamma=y_+-y_-,
\qquad
\phibase=\frac{y_++y_-}{\pi}.
\]
Here $y_++y_-$ is the unweighted Vieta sum of the two distinct carrier eigenvalues, not the trace on $E$.
\end{corollary}

\begin{proof}
By \Cref{eq:carrier-polynomial}, equivalently by the factorization
\[
6Y^2-Y-\mathbf 1=(2Y-\mathbf 1)(3Y+\mathbf 1),
\]
the two distinct carrier eigenvalues are exactly $y_+=1/2$ and $y_-=-1/3$. Their difference is
\[
y_+-y_-=\frac12+\frac13=\frac56=\gamma,
\]
while their sum is
\[
y_++y_-=\frac12-\frac13=\frac16.
\]
Since $\phibase=(1-\gamma)/\pi=1/(6\pi)$, the claim follows.
\end{proof}

\begin{proposition}[Dual-root generation of the hypercharge packet and admissible cusp set]
\label{prop:dual-root-generation}
Let
\[
y_+=\frac12,
\qquad
y_-=-\frac13.
\]
Then all one-family hypercharges are generated by the same two carrier roots:
\[
\begin{aligned}
Y(Q_L)&=y_++y_-,
&
Y(u^c)&=2y_-,
&
Y(d^c)&=-y_-,
\\
Y(L_L)&=-y_+,
&
Y(e^c)&=2y_+,
&
Y(\nu^c)&=0.
\end{aligned}
\]
Moreover the positive transport cusps satisfy
\[
\left\{1,\frac23,\frac13\right\}=\{2y_+,-2y_-,2(y_++y_-)\}.
\]
Thus the carrier packet and the transport grammar are generated by the same two-point root data.
\end{proposition}

\begin{proof}
On the minimal carrier one has
\[
Y=\operatorname{diag}(y_-,y_-,y_-,y_+,y_+)
\]
with respect to the split $E_3\oplus E_2$. The one-family packet $S^+$ carries the sectors
$Q_L,u^c,d^c,L_L,e^c,\nu^c$, and their hypercharges are the sums of the occupied carrier roots.
This gives the displayed formulas. Substituting $y_+=1/2$ and $y_-=-1/3$ recovers the canonical
values
\[
\frac16,\quad
-\frac23,\quad
\frac13,\quad
-\frac12,\quad
1,\quad
0.
\]
The admissible transport cusps are fixed in \TFPTxref{thm:exact-transport-closure} as
$\{1,\frac23,\frac13\}$, which is exactly the displayed root-generated set.
\end{proof}

\begin{corollary}[Family hypercharge variance and electroweak block trace]
Equip one carrier family with the uniform average
\[
\mathbb E_{S^+}(A):=\frac{1}{16}\tr_{S^+}A.
\]
Then
\[
\operatorname{Var}_{S^+}(Y)
:=
\mathbb E_{S^+}(Y^2)-\mathbb E_{S^+}(Y)^2
=
\frac{1}{16}\tr_{S^+}Y^2
=
\frac{5}{24}.
\]
Consequently,
\[
I_1^{\mathrm{EW}}=\operatorname{Var}_{S^+}(Y)\,\bone.
\]
\end{corollary}

\begin{proof}
The one-family trace identities give $\tr_{S^+}Y=0$ and $\tr_{S^+}Y^2=10/3$. Hence the centered second moment is
\[
\frac{1}{16}\tr_{S^+}Y^2=\frac{10/3}{16}=\frac{5}{24}.
\]
Using $I_1^{\mathrm{EW}}=\frac{5}{24}\,\bone$ from the electroweak block relation yields the second identity.
\end{proof}

\begin{corollary}[Factorial quadratic trace and two-defect factorization]
\label{cor:factorial-trace-delta2}
Let
\[
X:=6Y.
\]
Then the one-family quadratic trace is
\[
\tr_{S^+}(X^2)
=
36\,\tr_{S^+}(Y^2)
=
36\cdot\frac{10}{3}
=
120
=
5!.
\]
Equivalently, using the six hypercharge sectors with multiplicities
\[
(1,3,6,1,2,3)
\]
at
\[
X\in\{0,-4,1,6,-3,2\},
\]
the same trace is
\[
0^2+3\cdot 4^2+6\cdot 1^2+1\cdot 6^2+2\cdot 3^2+3\cdot 2^2=120.
\]
On the retained branch this gives the exact factorization
\[
\delta_2
=
2880\,\cthree^8
=
4!\,\tr_{S^+}(X^2)\,\cthree^8.
\]
\end{corollary}

\begin{proof}
By \TFPTxref{prop:split-weight-enumerator}, one has
\[
\tr_{S^+}(Y^2)=\frac{10}{3}.
\]
Since $X=6Y$, the first identity follows immediately. The sector-by-sector sum is the same quadratic trace written in the basis of admissible occupancy sectors from \TFPTxref{prop:family-hypercharge-compression}. Finally, the hard kernel normalization gives
\[
\delta_2=2880\,\cthree^8=24\cdot 120\,\cthree^8
=
4!\,\tr_{S^+}(X^2)\,\cthree^8.
\]
\end{proof}

\begin{remark}[Quadratic trace versus variance]
The value $120=5!$ is the unnormalized quadratic trace $\tr_{S^+}(X^2)$, not the normalized variance. The corresponding uniform variance is
\[
\operatorname{Var}_{S^+}(X)=36\,\operatorname{Var}_{S^+}(Y)=36\cdot\frac{5}{24}=\frac{15}{2}.
\]
\end{remark}

\begin{remark}[Generic abelian-index identity]
Define the generic occupancy
\[
\Omega_{\mathrm{occ}}:=16\,N_{\mathrm{fam}}.
\]
Then the abelian-index formula gives
\[
10\,\bone
=
\frac{40}{3}N_{\mathrm{fam}}+N_\Phi
=
\Omega_{\mathrm{occ}}\,\gamma+N_\Phi.
\]
At this stage no canonical specialization is inserted yet. The closed $41$ chain will be derived
only after the family and Higgs theorems fix $\Omega_{\mathrm{occ}}=48$ and $N_\Phi=1$.
\end{remark}

\begin{corollary}[Family integrality repair]
\label{cor:family-integrality-repair}
One carrier family contributes the weighted transport budget
\[
\dim S^+\cdot\gamma
=
16\cdot\frac56
=
\frac{40}{3}.
\]
Hence any integral global budget assembled from whole carrier families must satisfy
\[
N_{\mathrm{fam}}\cdot \frac{40}{3}\in\mathbb Z.
\]
The minimal positive repair is therefore
\[
N_{\mathrm{fam,min}}=3.
\]
For this minimal repair one gets
\[
\Omega_{\mathrm{occ}}=N_{\mathrm{fam,min}}\dim S^+=48,
\qquad
\Omega_{\mathrm{occ}}\gamma=40.
\]
The later family theorem realizes exactly this minimal arithmetic.
\end{corollary}

\begin{proof}
The first identity uses $\dim S^+=16$ and $\gamma=5/6$. Since $\gcd(40,3)=1$, integrality forces
$3\mid N_{\mathrm{fam}}$, so the smallest positive choice is $3$. Multiplying by $16$ and then by
$\gamma$ gives the displayed values.
\end{proof}

\subsection{Generic abelian index identity}

The abelian coefficient is first carried only in the generic form
\begin{equation}
\label{eq:41-chain-main}
10\,\bone
=
\Omega_{\mathrm{occ}}\gamma+N_\Phi.
\end{equation}
The arithmetic specialization to $41$ is deferred until the family multiplicity and the Higgs
sector have both been theorematically closed.


% ---- Families, occupancy, and compact Higgs index: source lines 3483-4124 ----
\section{Closure architecture A: families, Higgs, and local spectral scale}

This first closure-architecture block is the shortest route from the hard carrier packet to the
closed theorem body. It has three moves: primitive winding balance in the seam normal plane forces
the four-corner structure of the fundamental section, that winding theorem globalizes into the
family space $F$, the determinant line selects the unique light Higgs doublet together with the
compact seam phase, and the local spectral-scale channel generates both $\Mpl$ and the electroweak
vacuum from one common geometric response.

\subsection{Family monodromy and occupancy}

\begin{definition}[Family space]
Let $\overline X_f$ denote the compact seam quotient with four $\mathbb Z_4$ corner points
$p_1,\dots,p_4$, and define the punctured family section
\[
X_f^\circ:=\overline X_f\setminus\{p_1,p_2,p_3,p_4\}.
\]
Its rigid harmonic model is the Möbius class
\[
X_f^\circ\cong \mathbb P^1\setminus\mu_4.
\]
Set
\[
F:=H^1(X_f^\circ,\mathbb C).
\]
We use the distinguished cycle basis $(C_1,C_2,C_T)$ only as a concrete carrier-side basis of this family space; any later APS/Hodge realization is read as compatible analytic structure on the same topological family space.
\end{definition}
In this interpretation $F$ is neither an arbitrary flavor vector space nor a loose bookkeeping
trick; it is a cycle space attached to the seam data. The actual bridge logic is shown in
\Cref{fig:family-triangle}: the derived corner theorem fixes $N_{\mathrm{fam}}=3$, the
topological family step promotes that count to a family space, and the occupancy corollary then
feeds back into the defect channel.

\begin{figure}[t]
\centering
\scriptsize
\begin{tikzpicture}[
    >=Latex,
    hard/.style={draw=blue!60!black, fill=blue!6, rounded corners, align=center, text width=2.2cm, minimum height=1.0cm},
    closure/.style={draw=orange!75!black, fill=orange!8, rounded corners, align=center, text width=2.4cm, minimum height=1.0cm},
    bridge/.style={draw=green!50!black, fill=green!8, rounded corners, align=center, text width=2.25cm, minimum height=1.0cm}
]
\draw[draw=blue!60!black, fill=blue!4, rounded corners] (-5.5,-1.4) rectangle (-3.0,1.4);
\foreach \x/\y in {-5.1/1.0,-3.4/1.0,-5.1/-1.0,-3.4/-1.0} {
    \filldraw[blue!60!black] (\x,\y) circle (2.3pt);
    \node[blue!60!black, anchor=south] at (\x,\y) {\tiny $\frac34$};
}
\node[align=center] at (-4.25,0) {\textbf{corner compactification}\\four $\mathbb Z_4$ corners\\each contributes $\frac34$};

\node[hard] (nfam) at (-1.0,0) {$N_{\mathrm{fam}}$\\$=4\cdot\frac34=3$};

\coordinate (A) at (2.0,1.15);
\coordinate (B) at (0.7,-0.95);
\coordinate (C) at (3.3,-0.95);
\node[closure, text width=1.2cm, minimum height=0.8cm] at (A) {$C_1$};
\node[closure, text width=1.2cm, minimum height=0.8cm] at (B) {$C_2$};
\node[closure, text width=1.2cm, minimum height=0.8cm] at (C) {$C_T$};
\draw[->, thick, orange!75!black] ($(A)+(0.2,-0.3)$) .. controls (2.7,0.25) .. ($(C)+(-0.2,0.3)$);
\draw[->, thick, orange!75!black] ($(C)+(-0.38,-0.02)$) .. controls (2.0,-1.65) .. ($(B)+(0.38,-0.02)$);
\draw[->, thick, orange!75!black] ($(B)+(0.2,0.3)$) .. controls (1.3,0.2) .. ($(A)+(-0.2,-0.3)$);
\node[closure, text width=2.1cm, minimum height=0.9cm] at (2.0,-2.2) {topological family closure\\$F\cong\mathbb C^3$};

\node[closure] (occ) at (5.85,0.55) {$\Omega_{\mathrm{adm}}$\\$=3\times16=48$};
\node[bridge] (dtop) at (5.85,-1.15) {$\deltatop=\Omega_{\mathrm{adm}}\cthree^4$\\boundary-to-bulk bridge};

\draw[->, thick] (-3.0,0) -- (nfam);
\draw[->, thick] (nfam) -- (1.0,0.1);
\draw[->, thick] (nfam) -- (occ);
\draw[->, thick, orange!75!black] (2.0,-1.55) -- (2.0,-1.8);
\draw[->, thick, orange!75!black] (3.05,-1.9) -- (occ);
\draw[->, thick, green!50!black] (occ) -- (dtop);
\end{tikzpicture}
\caption{Carrier-to-family closure bridge. Corner compactification yields the punctured family
section $X_f^\circ$, the harmonic family space $\mathcal H_F\cong F\cong H^1(X_f^\circ,\mathbb C)$
gives the hard multiplicity $N_{\mathrm{fam}}=3$, and the resulting occupancy
$\Omega_{\mathrm{adm}}=48$ feeds back into the bridge reading
$\deltatop=\Omega_{\mathrm{adm}}\cthree^4$.}
\label{fig:family-triangle}
\end{figure}

The Hodge metric on $F$ turns the monodromy data into an honest unitary family system rather than a merely decorative triangle sketch.

\begin{theorem}[Four corners and the rigid family surface \texorpdfstring{$\mathbb P^1\setminus\mu_4$}{P1 minus mu4}]
\label{thm:topological-family-closure}
Let $\overline X_f$ be the compact seam quotient induced by the deck involution $\taudbl$ and its
spin lift on the oriented normal two-plane. Every admissible fixed corner contributes one
quarter-turn of the lifted spin rotation. Therefore
\[
N_{\mathrm{corner}}\cdot \frac{\pi}{2}=2\pi,
\qquad
N_{\mathrm{corner}}=4.
\]
Removing the four corner points gives a four-punctured sphere
\[
X_f^\circ:=\overline X_f\setminus\{p_1,p_2,p_3,p_4\}.
\]
The inherited finite symmetry group on the puncture set is a faithful dihedral group
$D_4\subset PGL_2(\mathbb C)$. Every faithful copy of $D_4$ in $PGL_2(\mathbb C)$ is M\"obius
conjugate to the standard action generated by
\[
r(z)=iz,
\qquad
s(z)=1/z.
\]
Up to M\"obius equivalence the unique $D_4$-invariant four-point configuration is therefore
\[
\mu_4:=\{z\in\mathbb C:\ z^4=1\}=\{1,-1,i,-i\}.
\]
Hence
\[
X_f^\circ\cong \mathbb P^1\setminus\mu_4,
\qquad
\chi(X_f^\circ)=-2,
\qquad
H_1(X_f^\circ,\mathbb Z)\cong \mathbb Z^3.
\]
Consequently
\[
F:=H^1(X_f^\circ,\mathbb C)\cong \mathbb C^3.
\]
\end{theorem}

\begin{proof}
The quarter-turn count gives four corner points. A sphere with four removed points is isomorphic to
$\mathbb P^1\setminus\{q_1,q_2,q_3,q_4\}$ for some ordered quadruple. The classification of finite
subgroups of $PGL_2(\mathbb C)$ implies that every faithful dihedral subgroup of order eight is
M\"obius conjugate to the standard $D_4$ action generated by $z\mapsto iz$ and $z\mapsto 1/z$.
Under that action the unique invariant orbit of cardinality four is the set of fourth roots of
unity, namely $\mu_4=\{1,-1,i,-i\}$. Therefore
\[
X_f^\circ\cong \mathbb P^1\setminus\mu_4.
\]
The Euler characteristic is $2-4=-2$, and for a genus-zero surface with four punctures one has
$\rank H_1=4-1=3$. Dualizing gives $F\cong\mathbb C^3$.
\end{proof}

\begin{theorem}[Three harmonic family modes on the four-punctured family surface]
\label{thm:harmonic-family-modes}
Let
\[
X_f^\circ=\mathbb P^1\setminus\mu_4
\]
with a complete cusp metric in its M\"obius class, and define the family Hilbert space by
\[
\mathcal H_F:=\ker \Delta^{(1)}_{L^2,\mathrm{APS}}(X_f^\circ).
\]
Then
\[
\dim\mathcal H_F=3.
\]
If the physical family multiplicity is defined by
\[
N_{\mathrm{fam}}:=\dim\mathcal H_F,
\]
then
\[
N_{\mathrm{fam}}=3,
\qquad
F\cong \mathcal H_F\cong \mathbb C^3.
\]
\end{theorem}

\begin{proof}
For a compact truncation $\overline X_{f,R}$ with four boundary circles,
\[
\chi(\overline X_{f,R})=2-4=-2.
\]
Hence
\[
\dim H^1(\overline X_{f,R},\partial\overline X_{f,R};\mathbb C)=3.
\]
The $L^2$ Hodge theorem for finite-area cusp surfaces identifies
\[
\ker \Delta^{(1)}_{L^2,\mathrm{APS}}(X_f^\circ)
\cong
H^1(\overline X_{f,R},\partial\overline X_{f,R};\mathbb C),
\]
so the dimension is $3$.
\end{proof}

\begin{remark}[How the family theorem is now used]
The family theorem is used at rigid conformal strength in the present version. Any orbifold APS
realization is read as compatible analytic structure on the same classification-rigid
four-punctured section, while the physical multiplicity itself is carried by the harmonic family
mode space $\mathcal H_F$.
\end{remark}

\begin{proposition}[Orbifold APS compatibility]
\label{prop:orbifold-aps-compatibility}
Under the additional analytical hypotheses of \TFPTxref{app:aps},
\[
\Ind^{\mathrm{orb,APS}}(D_{\mathrm{rel}}^+\otimes L_F)=3.
\]
\end{proposition}

\begin{corollary}[Boundary-to-bulk bridge corollary for the admissible occupancy]
\label{cor:boundary-bulk-occupancy}
\label{cor:derived-occupancy-identity}
On the canonical branch,
\[
\dim(S^+\otimes \mathcal H_F)=16\cdot 3=48=\Oadm.
\]
Hence the topological surcharge enters the present paper as a derived occupancy corollary:
\[
\deltatop
=
\cthree^4\,\dim(S^+\otimes \mathcal H_F)
=
\Oadm\,\cthree^4.
\]
\end{corollary}

The family subsection is therefore no longer split between a local count, an independent index
theorem, and a later occupancy reinterpretation. Those statements are compressed into one derived
corner theorem with one bridge corollary, while the holonomy proposition below records the
resulting unitary geometry on $F\cong \mathcal H_F$.

This theorem is the family statement referenced later by the admissibility theorems: the closure layer uses it as the explicit family component of the master closure theorem rather than as an external import.

\begin{proposition}[Rigid special-unitary family holonomy on the determinant-trivial branch]
\label{prop:family-holonomy-su3}
Let $L_F\to X_f^\circ$ be the flat Hermitian family bundle with fiber
$F\cong \mathcal H_F\cong H^1(X_f^\circ,\mathbb C)$. Assume that its determinant line is flat trivial,
\[
\det L_F \cong X_f^\circ\times \mathbb C,
\]
and that the chosen flat connection preserves this trivialization. Then
\[
\mathrm{Hol}(\nabla_F)\subset SU(3)_F.
\]
\end{proposition}

\begin{proof}
The Hermitian metric implies
\[
\mathrm{Hol}(\nabla_F)\subset U(3).
\]
The induced connection on $\det L_F$ has trivial holonomy by assumption, so for
every loop $\gamma$ one has
\[
\det \rho_F(\gamma)=1.
\]
Hence every holonomy matrix lies in $SU(3)$.
\end{proof}

The role of the present subsection is thus no longer only mnemonic. It displays both the local patch-deficit count and the global APS/Hodge realization before the later closure stack compresses them into the theorem-level output surface.

\subsection{Occupancy and boundary-to-bulk bridge}

With the family bridge installed, the matter space becomes $S^+\otimes \mathcal H_F$ and the
occupancy $\Oadm=48$ is already fixed by the harmonic family-mode theorem and its bridge
corollary. Here \emph{occupancy} means the admissible chiral-state count after family lifting; it
is not used as a primitive datum. The bridge point is not that family lifting creates a new
numerical surcharge. Rather, the inherited hard kernel value
\[
\deltatop=\frac{3}{256\pi^4}
\]
admits on the canonical family branch the occupancy reinterpretation $\deltatop=\Oadm\,\cthree^4$.
Likewise the hard second-order defect contribution is re-read as
\[
\delta_2=\frac54\,\deltatop^2=B\gamma\,\Oadm^2\cthree^8.
\]
In the present manuscript this occupancy corollary is used with a derived bridge reading:
it is the boundary-to-bulk step that later feeds both the electromagnetic closure and the
geometric channel. The coefficient $5/4$ is therefore not an extra numerical motif; it is the
minimal-carrier compression factor $B\gamma$ rewritten in defect language. At the same time, the
equality $\deltatop=\Oadm\cthree^4$ is not presented as an independent dynamical transgression
theorem in the ordinary field-theory sense.

\subsection{From bridge to geometric completion}

\begin{principle}[Closure split]
Once the carrier packet and family closure architecture are fixed, the focused main-text closure
problem splits into a geometric program $\Ageom$ and a transport program $\Atrans$. Record
algebra and prediction semantics are now tied directly to the admissible gap in the main theorem
stack, while E$_8$ scale grammar and horizon extrapolations are deferred to appendix-level
continuations.
\end{principle}

\[
D_{\mathrm{geo}}:=D_{\mathrm{ref}}+D_{\mathrm{rel}}=D.
\]
The two main-text programs are packaged as
\[
\Ageom
:=
\langle E_3\oplus E_2,\ D_{\mathrm{geo}},\ \mathcal B_{\mathrm{rel}},\ \wspin\rangle,
\qquad
\Atrans
:=
\langle U_6,\ T,\ D_y,\ \mathcal Y_y^{(\varepsilon)},\ \Ccyc\rangle.
\]
This split keeps the carrier theorem, the geometric reconstruction problem, and the record-level extensions from being flattened into one undifferentiated claim. In particular, the metric field is no longer treated here as an independent entry of $\Ageom$; it is reconstructed from the geometric Dirac anchor introduced next, while $D_{\mathrm{rel}}$ remains the comparison operator for spectral subtraction.

\subsection{Relative APS and superconnection setup}

The geometric completion is formulated on a graded bundle
\[
\mathcal E_{\mathrm{rel}}=\mathcal E_{\mathrm{rel}}^+\oplus\mathcal E_{\mathrm{rel}}^-,
\]
with the geometric Dirac anchor $D_{\mathrm{geo}}$, the declared reference operator
$D_{\mathrm{ref}}$, and APS-type boundary conditions on the seam $\Seam$; see
Appendix~\TFPTxref{app:aps}. The geometric package is
\[
\mathbb A_{\Seam}:=(D_{\mathrm{geo}},D_{\mathrm{rel}},A_{\Sigma}^{\mathrm{geo}}),
\qquad
A_{\Sigma}^{\mathrm{geo}}:=\mathcal B_{\mathrm{rel}}\oplus\nabla_{\chigeo}.
\]
Here $D_{\mathrm{geo}}$ is the geometric Dirac anchor, while $D_{\mathrm{rel}}$ is the relative comparison operator entering the spectral subtraction formulas. The bundle data $A_{\Sigma}^{\mathrm{geo}}$ reorganize the primitive seed package $A_{\Sigma}^{\mathrm{seed}}$ introduced earlier in \TFPTxref{eq:tfpt-master}.
\begin{theorem}[Geometric reconstruction from the Dirac anchor]
\label{thm:geometric-reconstruction-anchor}
\label{thm:geometric-reconstruction-from-Dgeo}
\label{thm:relative-geometric-reconstruction}
Assume
\[
Z(\mathcal A)\cong C^\infty(M),
\]
and the geometric anchor $D_{\mathrm{geo}}$ is of Dirac type with principal symbol
\[
\sigma(D_{\mathrm{geo}})(x,\xi)^2
=
g^{\mu\nu}(x)\xi_\mu\xi_\nu\,\mathbf 1.
\]
Then $D_{\mathrm{geo}}$ reconstructs the metric $g$, the spin structure, and the
Levi--Civita connection on the geometric branch. Inner fluctuations of the same anchored
spectral datum generate the gauge sector.
The operator $D_{\mathrm{rel}}$ is the reference-subtracted comparison operator on that
same branch and is not itself used as the local geometric anchor.
\end{theorem}

\begin{proof}[Proof sketch]
On the admissible branch the commutative center identifies the manifold algebra, while the
principal symbol of the Dirac anchor determines the quadratic form on cotangent vectors.
Standard Dirac-type reconstruction then recovers the spin bundle and Levi--Civita connection
from that same symbol data. The relative operator changes the comparison bookkeeping, not the
geometric source of the local symbol.
\end{proof}

\begin{theorem}[Retained branch from the admissibility selector]
\label{thm:retained-branch-from-Padm}
\label{thm:admissible-sector-on-reconstructed-geometry}
\label{thm:full-operator-reconstruction}
Let $D_{\mathrm{geo}}$ be the geometric anchor and let $\Padm$ be the retained
admissibility selector.
Then:
\begin{enumerate}[label=(\roman*), leftmargin=1.5em, topsep=3pt, itemsep=2pt]
\item $\Xbulk=\mathrm{Spec}\,Z(\mathcal A)$ and $\Seam=\mathrm{Fix}(\taudbl)$;
\item the local spectral density of $D_{\mathrm{geo}}$ determines the geometric response
      $\chigeo$;
\item $\Padm$ defines the retained admissible sector for observables, transport,
      and positivity statements on that branch.
\end{enumerate}
No claim is made that the compressed operator $\Padm D_{\mathrm{rel}}\Padm$
carries the local principal symbol of the geometry.
\end{theorem}

\begin{corollary}[No residual geometric datum beyond the geometric anchor]
The geometric package is reconstructed from $D_{\mathrm{geo}}$ and its local spectral density.
The selector $\Padm$ acts downstream on the retained sector and is not part of the
local geometric anchor.
\end{corollary}

\begin{definition}[Local spectral scale and canonical relative spectral residues]
\label{def:local-spectral-scale}
\label{def:canonical-residues}
Write the geometric response as a local spectral scale
\[
\chigeo=\Lambda e^\sigma,
\qquad
D_{\sigma,\mathrm{geo}}:=e^{-\sigma/2}D_{\mathrm{geo}}e^{-\sigma/2}.
\]
\[
\mathcal Z_{\mathrm{rel}}(s;\sigma)
:=
\operatorname{Tr}_{\mathrm{rel}}\!\left(
|D_{\sigma,\mathrm{geo}}+\mathcal B_{\mathrm{rel}}|^{-2s}
-|D_{\mathrm{ref}}|^{-2s}
\right)
+\frac{i\pi}{2}\Delta\eta_\Sigma.
\]
The canonical Einstein residue and order-zero residue are
\[
\mathcal E_D
:=
\operatorname*{Res}_{s=1}\mathcal Z_{\mathrm{rel}}(s;\sigma),
\qquad
\Gamma_D^{(4)}
:=
\operatorname{FP}_{s=0}\mathcal Z_{\mathrm{rel}}(s;\sigma).
\]
The gravitational branch is then read from the profile-free residue pair rather than from a freely
chosen spectral profile.
\end{definition}

\begin{corollary}[Constant-scale reduction of the geometric branch]
\label{cor:constant-scale-reduction}
If the local spectral scale is frozen to a constant $\sigma=\sigma_0$, then
$\chigeo=\chigeozero:=\Lambda e^{\sigma_0}$ and the present geometric branch reduces to the
constant-$\chigeo$ closure relations extracted from
$\Gamma_{\mathrm{grav}}:=-6\chigeo^2\mathcal E_D+\Gamma_D^{(4)}$:
\[
\Mpl^2=\frac{\chigeozero^2}{2\pi^2},
\qquad
\vgeo^2=Z_H\,\frac{\mu_\Phi^2(\chigeozero)}{\lambda_\Phi}.
\]
Once the unique stationary root of \TFPTxref{thm:absolute-spectral-planck-closure} is imposed, the
Einstein coefficient is equivalently the boundary-normalized readout
\[
\frac{\Mpl^2}{\lambda_\Sigma^2}=\frac{\rho_\star}{2\pi^2}.
\]
Thus the former constant-$\chigeo$ formulas are retained as the homogeneous branch of the local
spectral-scale formulation rather than as a competing package.
\end{corollary}


\subsection{Bosonic relative index and Higgs selection}

\begin{theorem}[Minimal determinant classes from the primitive charge generator]
\label{thm:determinant-class-higgs-selection}
\label{thm:determinant-clutching}
Let
\[
L_2:=\det E_2,
\qquad
L_3:=\det E_3
\]
over the compactified oriented normal sphere $S^2=D_N\cup_{S^1}D_S$. Then the primitive charge
generator
\[
X=-2P_3+3P_2
\]
together with center neutrality and seam-evenness forces the clutching map
\[
g_2(\varphi)=e^{i\varphi}
\]
for $L_2$, while center neutrality forces
\[
g_3(\varphi)=1
\]
for $L_3$. Therefore
\[
c_1(L_2)=\deg g_2=1,
\qquad
c_1(L_3)=\deg g_3=0.
\]
This pair is the unique minimal nonnegative admissible determinant class.
\end{theorem}

\begin{proof}
Complex line bundles on
\[
S^2=D_N\cup_{S^1}D_S
\]
are classified by their clutching maps along the equator:
\[
\pi_1(U(1))\cong \mathbb Z.
\]
Hence the first Chern class is exactly the winding degree of the transition function.

In the weak rank-two sector the primitive charge generator
\[
X=-2P_3+3P_2
\]
assigns the unique positive determinant winding to the seam-even weak block. The spin lift on the
oriented normal sphere acts there by the half-angle representation, so on the determinant line the
transition factor is
\[
g_2(\varphi)=e^{i\varphi},
\]
hence
\[
\deg g_2=1,
\qquad
c_1(L_2)=1.
\]

In the color sector admissibility imposes center neutrality. Therefore the determinant transition
function of the color block is null-homotopic and may be chosen constant:
\[
g_3(\varphi)=1.
\]
Consequently
\[
\deg g_3=0,
\qquad
c_1(L_3)=0.
\]

The full primitive seam twist is therefore exhausted in the weak determinant class. Any nonnegative
pair strictly smaller than $(1,0)$ would force vanishing weak determinant winding and hence remove
the seam-even weak line required by the primitive charge generator. Any larger nonnegative pair
would insert extra clutching not carried by the primitive generator and therefore violate
minimality. Thus the unique minimal nonnegative admissible determinant class is
\[
(c_1(L_2),c_1(L_3))=(1,0).
\]
\end{proof}

\begin{remark}[Energy interpretation on the retained branch]
The older asymptotic bosonic functional
\[
\mathcal E_{\mathrm{asym}}(\Phi_2,\Phi_3)
:=
\tau_2\bigl|\mathrm{wind}(\det\Phi_2)\bigr|
+\tau_3\bigl|\mathrm{wind}(\det\Phi_3)\bigr|
+m_3^2\dim\ker\mathcal B_{E_3}^+,
\]
with $\tau_3>\tau_2>0$, remains a useful physical intuition for why the retained branch is energetically preferred. In the present rewrite, however, the structural weight is carried by the determinant-class theorem above rather than by the cost functional itself.
\end{remark}

\begin{theorem}[Bosonic index and unique Higgs doublet from the minimal determinant class]
\label{thm:callias-higgs-selection}
\label{thm:compact-higgs-unique}
\label{thm:compact-determinant-higgs-index}
Compactify the oriented two-dimensional normal slice $N_\Seam$ to $S^2$. For $a\in\{2,3\}$
define
\[
\mathcal B_{E_a}:=\sigma^\mu\nabla^{(a)}_\mu+\Phi_a
\]
on the compactified normal sphere. Under the determinant-class data of
\Cref{thm:determinant-class-higgs-selection}, one has
\[
L_2\cong\mathcal O(1),
\qquad
L_3\cong\mathcal O,
\]
and hence
\[
H^0(S^2,\mathcal O(1))\cong\mathbb C^2,
\qquad
H^1(S^2,\mathcal O(1))=0,
\]
and the compact bosonic index satisfies
\[
\Ind(\mathcal B_{E_2}^+)=1,
\qquad
\Ind(\mathcal B_{E_3}^+)=0.
\]
Consequently the closed branch contains exactly one complex weak doublet and no light color triplet,
and the unique light seam-even bosonic zero mode lies in
\[
\Phi\in\Gamma(E_2)\cong (1,2)_{1/2},
\qquad
\taudbl^{\ast}\Phi=+\Phi.
\]
\end{theorem}

\begin{proof}
By \Cref{thm:determinant-class-higgs-selection}, the weak determinant line is the degree-one line
bundle
\[
L_2\cong\mathcal O(1)
\]
on
\[
S^2\cong\mathbb CP^1,
\]
whereas the color determinant class is trivial and
\[
L_3\cong\mathcal O.
\]
For $\mathbb CP^1$, Birkhoff--Grothendieck identifies every holomorphic line bundle by its degree,
and Riemann--Roch gives
\[
\chi(\mathcal O(1))=1+1=2.
\]
Serre duality yields
\[
H^1(S^2,\mathcal O(1))
\cong
H^0(S^2,\mathcal O(-3))^\vee
=0,
\]
so
\[
H^0(S^2,\mathcal O(1))\cong\mathbb C^2,
\qquad
H^1(S^2,\mathcal O(1))=0.
\]

The resulting two-dimensional zero-mode space is exactly one complex weak doublet, so the compact
bosonic index records it as
\[
\Ind(\mathcal B_{E_2}^+)=1.
\]

In the color sector one has $L_3\cong\mathcal O$, so no positive determinant degree is available.
Center neutrality forbids an unpaired seam-even color contribution, hence the positive and negative
triplet modes pair with zero net index:
\[
\Ind(\mathcal B_{E_3}^+)=0.
\]
Therefore the closed branch contains exactly one seam-even Higgs doublet and no light color
triplet, and the unique light bosonic zero mode lies in the weak block.
\end{proof}

\begin{corollary}[Bosonic rank is two]
\label{cor:bosonic-rank-two}
On the closed branch the positive carrier block has rank
\[
\dim E_+=2.
\]
Equivalently, the unique seam-even light bosonic zero mode fills one complex weak doublet and no
larger positive carrier block is retained.
\end{corollary}

\begin{proof}
By \TFPTxref{thm:callias-higgs-selection}, the light seam-even bosonic zero mode lies in the weak block
\[
\Phi\in\Gamma(E_2)\cong (1,2)_{1/2},
\]
and the compact index identifies its zero-mode space with
\[
H^0(S^2,\mathcal O(1))\cong\mathbb C^2.
\]
Hence the retained positive carrier block is exactly two-dimensional.
\end{proof}

This is the sharp form of the Higgs claim used here. The issue is not whether a Higgs doublet can
be written down, but whether a compact bosonic index can be shown to select it uniquely.
\begin{remark}[Callias shadow]
The older Callias formulation is retained only as the noncompact shadow of the same compact index
theorem and is discussed separately in the appendical analytic remarks. It no longer carries the
main proof burden.
\end{remark}

Once that uniqueness statement is fixed, the renormalizable carrier-compatible Yukawa couplings
are the standard ones,
\[
Q\,u^c\,\Phi,
\qquad
Q\,d^c\,\Phi^\dagger,
\qquad
L\,e^c\,\Phi^\dagger,
\qquad
L\,\nu^c\,\Phi,
\]
so the Higgs statement becomes a genuine structural bridge rather than a notation choice. The
formal carrier discharge from this Yukawa bridge is recorded later in
\TFPTxref{cor:bosonic-rank-two,thm:yukawa-forces-32,cor:yukawa-discharges-kpp}, once the closure-theorem Higgs block has
been stated in its final retained-branch form.


% ---- Discrete admissible branch and carrier-premise discharge: source lines 7479-8167 ----

The two-stage selector geometry is summarized in \Cref{fig:admissibility-stack}.

\begin{figure}[t]
\centering
\scriptsize
\begin{tikzpicture}[
    >=Latex,
    data/.style={draw=black!60, fill=gray!8, rounded corners, text width=2.55cm, minimum height=0.9cm, align=center},
    prim/.style={draw=blue!60!black, fill=blue!6, rounded corners, text width=2.55cm, minimum height=0.95cm, align=center},
    close/.style={draw=orange!75!black, fill=orange!8, rounded corners, text width=2.55cm, minimum height=0.95cm, align=center},
    result/.style={draw=green!50!black, fill=green!8, rounded corners, text width=3.55cm, minimum height=1.0cm, align=center},
    note/.style={font=\footnotesize\itshape, text=black!65, align=center}
]
\node[data] (seamdata) at (-3.3,1.95) {primitive datum\\$\iC,\ \APSdom,\ D_{\mathrm{rel}}$};
\node[data] (carrierdata) at (3.3,1.95) {carrier \& determinant closure\\$X,\ Z_c,\ Q_\Theta$};

\node[prim] (pseam) at (-4.55,0.15) {$P_{\Seam,+}$\\seam-even sector};
\node[prim] (pgap) at (-2.0,0.15) {$P_{\mathrm{gap}}$\\admissible gap};
\node[prim] (pprim) at (-3.3,-1.7) {$\Pprim$\\primitive selector};

\node[close] (psing) at (2.0,0.15) {$P_{\mathrm{sing}}$\\color singlets};
\node[close] (ptheta) at (4.55,0.15) {$P_\Theta$\\determinant sector};
\node[close] (pupgrade) at (3.3,-1.7) {$P_{\mathrm{sing}}P_\Theta$\\closure upgrade};

\node[result] (padm) at (0,-3.7) {$\Padm=\Pprim P_{\mathrm{sing}}P_\Theta$\\full admissibility projector};

\draw[->, thick] (seamdata) -- (pseam);
\draw[->, thick] (seamdata) -- (pgap);
\draw[->, thick] (carrierdata) -- (psing);
\draw[->, thick] (carrierdata) -- (ptheta);
\draw[->, thick] (pseam) -- (pprim);
\draw[->, thick] (pgap) -- (pprim);
\draw[->, thick] (psing) -- (pupgrade);
\draw[->, thick] (ptheta) -- (pupgrade);
\draw[->, thick] (pprim) -- (padm);
\draw[->, thick] (pupgrade) -- (padm);
\node[note] at (0,-5.0) {Primitive data determine $\Pprim$; carrier and determinant closure upgrade it to $\Padm$.};
\end{tikzpicture}
\caption{Two-stage upgrade from the primitive selector to the full admissibility projector. The
left lane is fixed by the primitive datum, while the right lane enters only after carrier and
determinant closure.}
\label{fig:admissibility-stack}
\end{figure}

\begin{theorem}[Calder\'on upgrade of the APS boundary domain]
\label{thm:apsdom-upgrade}
Let $\APSdomprov$ be the provisional APS realization of the geometric anchor used in
\TFPTxref{thm:well-posed-primitive-dynamics}, and let $D_b$ be the canonical doubling of
$D_{\mathrm{rel}}$ across $\Seam$. Let $C(D_b)$ denote the Calder\'on projector of $D_b$ on the
seam collar. Then there is a canonical boundary-elliptic upgrade
\[
\APSdomprov
\;\xrightarrow{\;C(D_b)\;}\;
\APSdom
\]
such that:
\begin{enumerate}[label=(\roman*), leftmargin=1.5em, topsep=3pt, itemsep=2pt]
    \item $\APSdom$ lies in the same admissible elliptic class as $\APSdomprov$ and preserves the
    self-adjoint Fredholm realization of the geometric anchor;
    \item the relative graph domain of $D_{\mathrm{rel}}$ is unchanged, so all primitive results
    that depended only on $\APSdomprov$ continue to hold with $\APSdom$;
    \item the family-pairing data $\APSfam$ are induced canonically on the doubled admissible
    branch and are not a separate primitive datum.
\end{enumerate}
\end{theorem}

\begin{proof}[Proof sketch]
The Calder\'on projector $C(D_b)$ is canonically associated with the doubled operator on the
seam collar and produces a boundary-elliptic projection compatible with the provisional APS
choice. Both $\APSdomprov$ and $\APSdom$ are admissible elliptic boundary conditions for the
same Dirac-type operator and therefore define equivalent self-adjoint Fredholm realizations of
the geometric anchor (boundary-elliptic equivalence preserves the Sobolev graph domain).
Consequently every statement made on the primitive scaffold under $\APSdomprov$ persists on the
closed branch under $\APSdom$. The induction of $\APSfam$ from canonical doubling is a standard
consequence of the cobordism structure of the doubled APS pair.
\end{proof}

\begin{proposition}[Master closure roadmap from canonical doubling]
\label{thm:master-closure}
\label{prop:master-closure-roadmap}
\label{ass:remaining-closure-inputs}
This proposition is a roadmap statement: it lists the closure conclusions that the individual
theorems of this section establish from $\Tbound$, the derived closed datum, the carrier theorem, and the
canonical doubling. It is not used as an input to those proofs and is recovered as a master
corollary at the end of the section. Let $X_b$ be the canonical double across $\Seam$ and write
\[
D_b=
\begin{pmatrix}
0&D_{\mathrm{rel}}^-\\
D_{\mathrm{rel}}^+&0
\end{pmatrix},
\qquad
Q_{\mathrm{tot}}:=Q_{\mathrm{adm}}+Q_{\mathrm{geo}},
\qquad
\Delta_{\mathrm{tot}}:=\{Q_{\mathrm{tot}},Q_{\mathrm{tot}}^\dagger\}.
\]
Then on the closed branch all of the following are derived from the same datum (each item is the
output of a downstream theorem in this section, not an extra hypothesis here):
\begin{enumerate}[label=(\arabic*), leftmargin=1.5em, topsep=3pt, itemsep=2pt]
    \item the Calder\'on projector of $D_b$ canonically upgrades the provisional APS choice
    $\APSdomprov$ to $\APSdom$ (cf.\ \TFPTxref{thm:apsdom-upgrade});
    \item the family pairing data are induced canonically, so no separate $\APSfam$ datum remains;
    \item the upgrade theorem \TFPTxref{thm:padm-upgrade} produces
    \[
    \Padm:=\Pi_{\ker \Delta_{\mathrm{adm}}}
    =
    \Pprim\,P_{\mathrm{sing}}\,P_\Theta,
    \]
    and the doubled total closure package preserves $\mathrm{Ran}(\Padm)$;
    \item
    \[
    \chi(X_f^\circ)=-2,
    \qquad
    F\simeq H^1(X_f^\circ,\mathbb C)\simeq \mathbb C^3,
    \qquad
    N_{\mathrm{fam}}=3;
    \]
    \item the harmonic determinant line is parallel and therefore
    \[
    \mathrm{Hol}(\nabla_F)\subset SU(3)_F;
    \]
    \item admissible bosonic, fermionic, and geometric reflection positivity hold on $\mathrm{Ran}(\Padm)$;
    \item the physical geometric sector takes the form
    \[
    \mathcal H_{\mathrm{phys}}
    \cong
    \mathcal H_2\otimes\mathcal H_{\mathrm{sc}}\otimes\mathcal H_{\mathrm{matt}},
    \]
    with the scalaron carried by the positive $(\chi_{\mathrm{geo}},R^2)$ block;
    \item $\Padm\Hrel\Padm$ is lower bounded and coercive.
\end{enumerate}
\end{proposition}

\begin{definition}[Relative generating functional]
\label{def:relative-generating-functional}
The relative Euclidean generating functional is normalized as
\[
Z_{\mathrm{rel}}[J,\eta,\bar\eta]
:=
\frac{Z_{\mathrm{adm}}[J,\eta,\bar\eta]}{Z_{\mathrm{ref}}[0,0,0]}.
\]
The reference sector is therefore a fixed normalizing denominator rather than a second fluctuating measure that would mix with admissible positivity.
\end{definition}

\begin{lemma}[Admissible reflection positivity]
\label{lem:admissible-reflection-positivity}
Assume
\[
[\Theta,P_\Theta]=0,
\qquad
\Theta\Padm=\Padm\Theta,
\]
the admissible bosonic kernel satisfies the Bernstein-positivity remark above and is
reflection-stable on the admissible sector, and the
admissible fermionic CAR kernel is reflection positive on that same sector. Then every
admissible observable $O$ supported in $\tau>0$ satisfies
\[
\langle \Theta O,O\rangle_{\mathrm{rel}}
=
\frac{1}{Z_{\mathrm{ref}}[0,0,0]}
\langle \Theta O,O\rangle_{\mathrm{adm}}
\ge 0.
\]
\end{lemma}

\begin{proof}
By definition,
\[
\langle \Theta O,O\rangle_{\mathrm{rel}}
=
\frac{1}{Z_{\mathrm{ref}}[0,0,0]}\,
\langle \Theta O,O\rangle_{\mathrm{adm}}.
\]
The denominator is a fixed positive normalizing constant, so the sign is determined entirely by the
admissible numerator.

Split $O$ into its bosonic and fermionic parts on the retained admissible sector. For the bosonic
part, the Bernstein representation of the admissible kernel writes the Euclidean covariance as a
positive superposition of reflection-stable heat kernels. Reflection across $\tau=0$ therefore sends
every positive-time bosonic observable to a positive quadratic form, so the bosonic contribution to
\[
\langle \Theta O,O\rangle_{\mathrm{adm}}
\]
is nonnegative.

For the fermionic part, the CAR reflection-positivity hypothesis on the admissible subspace implies
that the reflected two-point kernel is positive in the standard fermionic sense; equivalently, the
corresponding Pfaffian / determinant form on positive-time test spinors is nonnegative. Because the
bosonic and fermionic admissible factors commute inside the normalized expectation, their product
remains nonnegative. Dividing by the fixed reference normalization preserves that sign. Hence
\[
\langle \Theta O,O\rangle_{\mathrm{rel}}\ge 0.
\]
\end{proof}

\begin{theorem}[Operative admissibility selector]
Under \TFPTxref{def:full-admissibility-complex,thm:padm-upgrade,prop:master-closure-roadmap}, the
operator $\Padm$ is idempotent on the retained branch and realizes the admissibility predicate
on all sectors used by the closed output statements:
\[
\mathsf{Adm}(X)=1
\quad\Longleftrightarrow\quad
\Padm X = X.
\]
\end{theorem}

\begin{proof}[Proof sketch]
By \TFPTxref{thm:padm-upgrade}, $\Padm=\Pprim P_{\mathrm{sing}}P_\Theta$ is the orthogonal projector
onto the harmonic admissible sector of the full complex and is therefore idempotent by
construction. \TFPTxref{cor:factorized-admissibility-selector} recovers the explicit factorized
selector. The predicate and operator languages agree once the positive transport sector is fixed.
\end{proof}

\begin{theorem}[Topological family closure on the admissible branch]
By \TFPTxref{thm:master-closure,thm:topological-family-closure,cor:boundary-bulk-occupancy},
\[
\chi(X_f^\circ)=-2,
\qquad
F\cong\mathbb C^3,
\qquad
N_{\mathrm{fam}}=3.
\]
Consequently
\[
\Oadm=N_{\mathrm{fam}}\dim S^+=3\times 16=48,
\qquad
\deltatop=\Oadm\,\cthree^4.
\]
\end{theorem}

\begin{proof}[Proof sketch]
This theorem is now a direct consequence of the master closure theorem from Section~6. The
topological family closure and the orbifold corner count have already been compressed there
into one statement, so the present layer merely records the corresponding admissibility
consequence.
\end{proof}

\begin{theorem}[Bosonic index and local spectral-scale closure]
On the admissible branch,
\[
\Ind_{\mathrm{rel}}^{\Padm}(\mathcal B_{E_2}^+)=1,
\qquad
\Ind_{\mathrm{rel}}^{\Padm}(\mathcal B_{E_3}^+)=0,
\]
so the unique light seam-even bosonic zero mode is
\[
\Phi\in\Gamma(E_2)\cong(1,2)_{1/2},
\qquad
\taudbl^{\ast}\Phi=+\Phi.
\]
Moreover the same closure layer yields
\[
\begin{aligned}
\frac{\Mpl^2}{\lambda_\Sigma^2}&=\frac{\rho_\star}{2\pi^2},
&
G_N\lambda_\Sigma^2&=\frac{\pi}{4\rho_\star},
\\
\frac{\vgeo}{\Mpl}&=\gcar\,\beta_{\mathrm{rad}}^2
\exp\!\left[-\frac{\alpha^{-1}(0)+\delta_{\mathrm{ph}}}{5}\right],
\\
G_N\vgeo^2&=\frac{1}{8\pi}\gcar^2\beta_{\mathrm{rad}}^4
\exp\!\left[-\frac{2(\alpha^{-1}(0)+\delta_{\mathrm{ph}})}{5}\right].
\end{aligned}
\]
Here $\vgeo$ is the seam-even UV source scale. The physical electroweak benchmark is the residue
matched quantity $v_{\mathrm{phys}}$ of \TFPTxref{thm:ew-geometric-to-physical-matching}.
\end{theorem}

\begin{proof}[Proof sketch]
The bosonic sector is treated as a genuine Callias / Dirac / Dolbeault-type index problem rather
than a scalar counting exercise. Relative heat-kernel closure together with
\TFPTxref{prop:canonical-bosonic-normalization,lem:two-sheet-zero-mode-normalization,thm:absolute-spectral-planck-closure}
produces the boundary-normalized Einstein branch, the canonical Higgs normalization, and the
seam-even zero-mode formula for $\vgeo$.
\end{proof}

\begin{definition}[Primitive Yukawa generator]
\label{def:primitive-yukawa-generator}
Let
\[
\mathcal I_Y
:=
\bigoplus_{p,q,r,t}
\operatorname{Hom}_{G_{\mathrm{car}}}
\Bigl(
(\Lambda^p E_-\otimes \Lambda^q E_+)
\otimes
(\Lambda^r E_-\otimes \Lambda^t E_+)
\otimes E_+,\,
\mathbb C
\Bigr)
\]
be the graded algebra of carrier-invariant Yukawa trilinears. A homogeneous nonzero invariant is
called primitive if it is indecomposable, equivalently if it does not lie in the ideal generated by
positive-degree invariants of strictly lower carrier degree.
\end{definition}

\begin{lemma}[Lowest primitive Yukawa generator on the closed branch]
\label{lem:lowest-primitive-yukawa-generator}
Assume $\dim E_+=2$ and let $\mathbb Y_{\mathrm{br}}$ be the branch Yukawa tensor of
\TFPTxref{thm:exact-transport-closure,thm:rigid-family-local-system}. Then the primitive generator of
$\mathcal I_Y$ with two nontrivial fermionic legs is unique up to exchange of the two fermionic
legs and belongs to
\[
\operatorname{Hom}_{G_{\mathrm{car}}}
\bigl((E_-\otimes E_+)\otimes \Lambda^2 E_-\otimes E_+,\,\mathbb C\bigr).
\]
\end{lemma}

\begin{proof}
The positive-block determinant neutrality gives $q+t+1=2$, hence $q+t=1$. Up to exchanging the two
fermionic legs, $(q,t)=(1,0)$. Evenness of both fermionic legs then forces $p$ odd and $r$ even.
If $p\ge 3$, the corresponding invariant factors through an extra negative spectator wedge and
therefore lies in the ideal generated by lower carrier degree. Likewise, if $r\ge 4$, the
invariant factors through a spectator bivector insertion and is again decomposable. Primitive
indecomposability therefore forces $p=1$ and $r=2$.
\end{proof}

\begin{theorem}[Carrier rigidity from the primitive component of the branch Yukawa tensor]
\label{thm:yukawa-forces-32}
Let
\[
E=E_-\oplus E_+
\]
be the two-factor carrier and let
\[
G_{\mathrm{car}}:=S\bigl(U(E_-)\times U(E_+)\bigr),
\qquad
S^+=\Lambda^{\mathrm{even}}E.
\]
Let $\mathbb Y_{\mathrm{br}}$ denote the branch Yukawa tensor induced by the canonical transport
kernel and rigid family holonomy of
\TFPTxref{thm:exact-transport-closure,thm:rigid-family-local-system}. Decompose its carrier-homogeneous
components as trilinear maps of the form
\[
B_{p,q;r,t}:
(\Lambda^p E_-\otimes \Lambda^q E_+)
\otimes
(\Lambda^r E_-\otimes \Lambda^t E_+)
\otimes E_+
\longrightarrow
\mathbb C.
\]
Let the primitive component of $\mathbb Y_{\mathrm{br}}$ be the unique indecomposable generator of
\Cref{def:primitive-yukawa-generator,lem:lowest-primitive-yukawa-generator} with two nontrivial
fermionic legs. Then, up to exchanging the two fermionic legs, that primitive component is
necessarily of type
\[
(E_-\otimes E_+) \otimes \Lambda^2 E_- \otimes E_+
\longrightarrow
\Lambda^3 E_-\otimes \Lambda^2 E_+\cong\mathbb C,
\]
and therefore
\[
\dim E_+=2,
\qquad
\dim E_-=3.
\]
Consequently
\[
\dim \Lambda^2 E_+=1,
\qquad
\Lambda^2 E_- \cong E_-^\vee,
\qquad
\dim \Lambda^{\mathrm{even}}E=16.
\]
Hence the rigid carrier is
\[
E=E_3\oplus E_2,
\qquad
X=-2P_3+3P_2,
\qquad
Y=-\frac13 P_3+\frac12 P_2.
\]
\end{theorem}

\begin{proof}
By \TFPTxref{cor:bosonic-rank-two}, the compact Higgs index has already fixed
\[
\dim E_+=2.
\]
Write the primitive nonzero Yukawa component with two nontrivial fermionic legs as
\[
B_{p,q;r,t}:
(\Lambda^p E_-\otimes \Lambda^q E_+)
\otimes
(\Lambda^r E_-\otimes \Lambda^t E_+)
\otimes E_+
\longrightarrow
\mathbb C
\]
and note that \Cref{lem:lowest-primitive-yukawa-generator} forces primitiveness in the invariant-ring
sense to have
\[
(q,t)=(1,0),
\qquad
(p,r)=(1,2)
\]
up to exchanging the two fermionic legs. Thus the primitive carrier component is exactly
\[
(E_-\otimes E_+) \otimes \Lambda^2 E_- \otimes E_+.
\]
Its negative factor contracts canonically as
\[
E_- \otimes \Lambda^2 E_- \longrightarrow \Lambda^3 E_-,
\]
so scalarity forces $\Lambda^3 E_-$ to be the top exterior power. Hence
\[
\dim E_-=3.
\]
With
\[
(\dim E_-,\dim E_+)=(3,2),
\]
one obtains the carrier-signature consequences
\[
\Lambda^2 E_- \cong E_-^\vee,
\qquad
\dim\Lambda^2 E_+=1,
\]
and therefore
\[
\dim \Lambda^{\mathrm{even}}(E_3\oplus E_2)=2^{5-1}=16.
\]
Finally the unimodular primitive two-point generator is fixed by
\[
3q_-+2q_+=0,
\qquad
q_-<0<q_+,
\qquad
\gcd(q_-,q_+)=1,
\]
so $(q_-,q_+) = (-2,3)$ and hence
\[
X=-2P_3+3P_2,
\qquad
Y=\frac{X}{6}=-\frac13 P_3+\frac12 P_2.
\]
\end{proof}

\begin{lemma}[Uniqueness of the cubic weak invariant]
\label{lem:unique-cubic-weak-invariant}
Let $E=E_-\oplus E_+$ be a two-factor carrier on the retained branch and let
\[
\Phi\in E_+
\]
be the unique seam-even light boson selected by the compact bosonic index. Then
\[
\dim \operatorname{Hom}_{G_{\mathrm{car}}}\bigl((E_-\otimes E_+)\otimes \Lambda^2 E_-\otimes E_+,\,\mathbb C\bigr)=1
\]
if and only if
\[
(\dim E_-,\dim E_+)=(3,2).
\]
In particular, on the retained branch the up-type cubic invariant is not an extra choice but the
unique admissible cubic invariant.
\end{lemma}

\begin{proof}
By Schur functor decomposition,
\[
\begin{aligned}
\operatorname{Hom}_{G_{\mathrm{car}}}\bigl((E_-\otimes E_+)\otimes \Lambda^2 E_-\otimes E_+,\mathbb C\bigr)
&\cong
\operatorname{Hom}_{SU(E_-)}(E_-\otimes \Lambda^2E_-,\mathbb C)
\\
&\qquad\otimes
\operatorname{Hom}_{SU(E_+)}(E_+\otimes E_+,\mathbb C).
\end{aligned}
\]
The first factor is nonzero iff $\Lambda^2E_-\cong E_-^\vee$, hence iff $\dim E_-=3$. The second
factor is nonzero iff the alternating contraction on two fundamentals is one-dimensional, hence iff
$\dim E_+=2$. In that case both factors are one-dimensional, so the total invariant space is
one-dimensional.
\end{proof}

\begin{corollary}[Carrier-signature premises are discharged internally]
\label{cor:yukawa-discharges-kpp}
Combine \TFPTxref{cor:bosonic-rank-two,thm:yukawa-forces-32,lem:unique-cubic-weak-invariant} with the
one-family decomposition of
$S^+=\Lambda^{\mathrm{even}}(E_3\oplus E_2)$ and with the unique seam-even Higgs theorem. Then the
weak antisymmetric line, the color-dual antisymmetric sector, and the retained light-packet
identification are no longer primitive inputs once the uniqueness of the retained cubic invariant
has been proved on the same branch.
\end{corollary}

\begin{figure}[t]
\centering
\scriptsize
\resizebox{\linewidth}{!}{%
\begin{tikzpicture}[
    >=Latex,
    source/.style={draw=blue!60!black, fill=blue!6, rounded corners, text width=2.75cm, minimum height=0.95cm, align=center},
    force/.style={draw=orange!75!black, fill=orange!8, rounded corners, text width=2.9cm, minimum height=0.95cm, align=center},
    result/.style={draw=green!50!black, fill=green!8, rounded corners, text width=2.9cm, minimum height=0.95cm, align=center},
    note/.style={font=\footnotesize\itshape, text=black!60, align=center}
]
\node[source] (fam) at (-5.2,0.95) {$X_f^\circ,\ N_{\mathrm{fam}}=3$\\family closure};
\node[source] (higgs) at (-1.7,0.95) {$\Ind(\mathcal B_{E_2}^+)=1$\\$\dim E_+=2$};
\node[force] (yuk) at (1.8,0.95) {primitive indecomposable\\Yukawa generator};
\node[result] (car) at (5.3,0.95) {$E_3\oplus E_2,\ g=5,$\\$S^+=\Lambda^{\mathrm{even}}E$};

\node[result] (cons) at (1.8,-1.2) {derived check\\$\Oadm=48=4\cdot 12$};

\draw[->, thick] (fam) -- (higgs);
\draw[->, thick] (higgs) -- (yuk);
\draw[->, thick] (yuk) -- (car);
\draw[->, thick] (car) -- (cons);
\draw[->, thick, dashed] (fam) -- (cons);

\node[note] at (0,-2.35) {Carrier closure runs through Higgs rank and branch Yukawa rigidity; the
old occupancy formula is only a downstream consistency identity.};
\end{tikzpicture}%
}
\caption{Carrier closure after the proof-tree cut. The retained branch first fixes $\dim E_+=2$
through the compact Higgs index, then $\dim E_-=3$ through the primitive indecomposable generator of
the actual branch Yukawa tensor; only afterwards is $\Oadm=48=4\cdot12$ read as a consistency check.}
\label{fig:carrier-closure-cut}
\end{figure}

\begin{corollary}[Hard holonomy form of the Yukawa matrices]
On the admissible branch, each fermion sector obeys the canonical factorization
\[
Y_f=D_{L,f}U_f^\star D_{R,f},
\qquad
U_f^\star\in SU(3)_F,
\]
and, in the family basis fixed by \TFPTxref{thm:exact-transport-closure},
\[
(Y_f)_{ij}
=
\lambdaY^{L^L_{f,i}+L^R_{f,j}}
\left\langle e_i,\mathrm{Hol}_{\Gamma_{ij}^{\min}}(\nabla_F^\star)e_j\right\rangle.
\]
Any factorization by $U_f^\star$ is therefore a consequence of the hard holonomy closure, not an
independent flavor input. The diagonal path-length rows \TFPTxref{eq:path-length-yukawa} are the
singular-value reduction of this matrix statement.
\end{corollary}

\begin{proof}[Proof sketch]
The finite $C_6$ transport geometry fixes the graph-distance weights, while the rigid
$D_4$-equivariant family connection fixes the holonomy factors on the canonical path frame. The
diagonal ladder formulas arise once left and right admissible words are paired on the same
singular-value channel, so the matrix factorization is recovered as a corollary of one closed
holonomy statement rather than as a separate ansatz.
\end{proof}

\begin{definition}[Center-neutral admissible algebra]
Let
\[
\Asing
:=
\overline{\mathrm{span}}\{A_w:\ z(w)=1,\ \varepsilon(w)=+1,\ \Padm A_w=A_w\}.
\]
\end{definition}

\begin{theorem}[Hadronic admissibility closure]
On the admissible branch, the physical hadronic Hilbert space is
\[
\Hhad:=\overline{\Asing|\Omega\rangle}.
\]
All physical hadronic states lie in $\Hhad$, while color-nonsinglet free-word sectors are removed by $\Padm$.
\end{theorem}

\begin{proof}
Center neutrality and seam-evenness are already the carrier-side admissibility rules for physical hadrons. Passing to the norm closure of the singlet admissible algebra acting on the vacuum gives the completed hadronic space, and no state with $z(w)\neq 1$ survives the selector.
\end{proof}

\begin{theorem}[Anomaly cancellation and seam inflow]
On the admissible branch,
\[
I_6^{\mathrm{tot}}=0,
\qquad
\delta_\Lambda S_{\mathrm{bulk}}+\delta_\Lambda S_{\eta,\Seam}=0.
\]
\end{theorem}

\begin{proof}[Proof sketch]
The one-family packet $S^+$ is exactly the anomaly-free Standard Model family including $\nu^c$. The remaining seam contribution is carried by the relative $\eta$ term, so gauge variation cancels between bulk and seam sectors rather than being left as an uncancelled boundary defect.
\end{proof}

\begin{lemma}[Topological sector positivity]
\label{lem:topological-sector-positivity}
On the admissible branch every topological sector weight satisfies
\[
Z_Q\ge 0.
\]
\end{lemma}

\begin{proof}
Apply \TFPTxref{lem:retained-gamma5-hermiticity}. The nonzero spectrum of each $\mathcal D_f$ occurs in
paired sets $\pm i\lambda_n$, while the branch mass eigenvalues are strictly positive by
\TFPTxref{thm:sheet-cp-protection}. Therefore each sector determinant is nonnegative, and so is the
full sector weight.
\end{proof}

\begin{lemma}[Nontrivial topological support on the closed branch]
\label{lem:nontrivial-topological-support}
On the closed branch the primitive topological sector has strictly positive weight:
\[
Z_1=Z_{-1}>0.
\]
\end{lemma}

\begin{proof}
By \TFPTxref{thm:boundary-winding-control}, the canonical branch has unit seam winding, so the primitive
topological class $Q=1$ is nonempty on the admissible branch. Choose an admissible representative
$\phi_1$ in that class. By \TFPTxref{thm:sheet-cp-protection}, the same branch carries a genuine compact
determinant angle with primitive winding number one, so $\phi_1$ has finite admissible action and
belongs to the closed branch sector $Q(\phi_1)=1$. Let $\mathcal U_1$ be a sufficiently small open
neighborhood of $\phi_1$ inside that sector. The local Boltzmann density is strictly positive on
$\mathcal U_1$, and \Cref{lem:topological-sector-positivity} gives nonnegative sector weights.
Therefore
\[
Z_1
\ge
\int_{\mathcal U_1} e^{-S_{\mathrm{adm}}[\phi]}\,d\mu(\phi)
>
0.
\]
The antiunitary sheet symmetry exchanges $Q=1$ and $Q=-1$, hence $Z_{-1}=Z_1>0$.
\end{proof}

\begin{theorem}[Unconditional strong-CP closure on the admissible branch]
\label{thm:strong-cp-closure}
\label{thm:strong-cp-closure-positive}
On the admissible branch one has
\[
\arg\det M_u=\arg\det M_d=0,
\qquad
\bar\theta=0.
\]
Moreover the free energy
\[
F(\theta):=-V^{-1}\log Z_{\mathrm{adm}}(\theta)
\]
has its unique minimum at
\[
\theta=0\!\!\mod 2\pi.
\]
Weak CP may still survive through
\[
V_{\mathrm{CKM}}=U_{u,L}^\dagger U_{d,L}.
\]
\end{theorem}

\begin{proof}
\TFPTxref{thm:sheet-cp-protection} gives
\[
\arg\det M_u=\arg\det M_d=0.
\]
By \Cref{lem:topological-sector-positivity}, the admissible partition function admits a sector
decomposition
\[
Z_{\mathrm{adm}}(\theta)=\sum_{Q\in\mathbb Z} Z_Q e^{iQ\theta},
\qquad
Z_Q\ge 0.
\]
The antiunitary sheet symmetry of \TFPTxref{thm:sheet-cp-protection} gives
\[
\mathcal C_\Sigma\,Q\,\mathcal C_\Sigma^{-1}=-Q,
\]
so the sector weights satisfy $Z_Q=Z_{-Q}$. Hence
\[
\begin{aligned}
Z_{\mathrm{adm}}(\theta)
&=Z_0+2\sum_{Q>0}Z_Q\cos(Q\theta)\\
&\le Z_0+2\sum_{Q>0}Z_Q
=Z_{\mathrm{adm}}(0).
\end{aligned}
\]
By \Cref{lem:nontrivial-topological-support}, one has $Z_1>0$. Therefore for every
$\theta\not\equiv 0\mod 2\pi$,
\[
Z_{\mathrm{adm}}(\theta)
\le
Z_{\mathrm{adm}}(0)-2Z_1\bigl(1-\cos\theta\bigr)
<
Z_{\mathrm{adm}}(0),
\]
which implies
\[
F(\theta)>F(0)
\qquad
\text{for every }\theta\not\equiv 0\mod 2\pi.
\]
Thus the unique minimum is at $\theta=0$, and therefore $\bar\theta=0$.
\end{proof}

% ---- Joint discrete admissible sector excerpt: source lines 9158-9639 ----
\begin{theorem}[Joint discrete admissible sector from the boundary primitive kernel]
\label{thm:joint-discrete-admissible-sector}
\label{thm:discrete-sector-from-primitive-kernel}
For every admissible boundary primitive kernel $\Tker^\partial$, the compatible discrete datum
\[
\mathfrak D_{\mathrm{disc}}(\Tker^\partial)
\]
is a singleton. Its unique element is
\[
d^\star_{\mathrm{disc}}
:=
\Bigl(
E_3\oplus E_2,\,
Y,\,
S^+,\,
X_f^\circ,\,
[\nabla_F^\star],\,
c_1(L_2),\,
c_1(L_3),\,
N_{\mathrm{fam}},\,
\theta_{\mathrm{eff}}
\Bigr)
\]
with
\[
X_f^\circ\cong\mathbb P^1\setminus\mu_4,
\qquad
N_{\mathrm{fam}}=3,
\qquad
(g,b,s)=(5,3,2),
\qquad
\dim S^+=16,
\]
\[
\begin{aligned}
S^+&=\Lambda^{\mathrm{even}}(E_3\oplus E_2),\\
Y&=\mathrm{diag}\!\left(-\tfrac13,-\tfrac13,-\tfrac13,\tfrac12,\tfrac12\right),\\
(c_1(L_2),c_1(L_3))&=(1,0),\qquad
\theta_{\mathrm{eff}}=0.
\end{aligned}
\]
Equivalently, carrier data, family data, determinant data, and the effective strong angle are fixed
jointly from the same boundary primitive kernel.
\end{theorem}

\begin{proof}
We prove the discrete blocks in the only order that uses no downstream uniqueness claim.

\textbf{Step 1: primitive winding.}
By \TFPTxref{thm:primitive-seam-generator}, the primitive seam class is the positive generator
\[
[u_\Sigma]=1\in K^1(S^1)\cong\mathbb Z,
\]
and its associated spectral flow is
\[
\SF(U_\Sigma)=1.
\]
Hence the primitive branch carries exactly one elementary seam winding. Write this winding number as
\[
n:=\SF(U_\Sigma)=1.
\]

\textbf{Step 2: family geometry and family multiplicity.}
The lifted spin action on the oriented normal two-plane closes only after $4\pi$, so one primitive
winding is realized by four quarter-turn corners. Therefore the compact family quotient carries
exactly four $\mathbb Z_4$ orbifold points, and removing them produces
\[
X_f^\circ\cong\mathbb{CP}^1\setminus\{q_1,q_2,q_3,q_4\}.
\]
The lifted seam symmetry acts by the symmetries of a square, so the marked configuration admits a
faithful $D_4$ action. By the classification of finite subgroups of $PGL_2(\mathbb C)$, every
faithful $D_4$ action on four marked points of $\mathbb{CP}^1$ is M\"obius conjugate to the square
configuration
\[
\{1,-1,i,-i\}.
\]
Hence
\[
X_f^\circ\cong\mathbb P^1\setminus\mu_4.
\]
The corner deficit count is then
\[
N_{\mathrm{fam}}
:=
\sum_{j=1}^{4}\left(1-\frac14\right)
=
4\cdot\frac34
=
3.
\]
Because $\chi(X_f^\circ)=2-4=-2$, one also has
\[
H_1(X_f^\circ,\mathbb Z)\cong\mathbb Z^3,
\qquad
H^1(X_f^\circ,\mathbb C)\cong\mathbb C^3.
\]
Thus the family multiplicity and the family geometry are fixed already at the discrete level.

\textbf{Step 3: carrier closure.}
By \TFPTxref{cor:bosonic-rank-two}, the compact Higgs index has already fixed
\[
\dim E_+=2.
\]
By \TFPTxref{thm:yukawa-forces-32}, the primitive indecomposable generator of the actual branch Yukawa tensor
then forces
\[
\dim E_-=3.
\]
Hence
\[
\begin{aligned}
(b,s)&=(3,2),\\
g&=b+s=5,\\
\dim S^+&=2^{g-1}=16.
\end{aligned}
\]

\textbf{Derived consistency check.}
With $N_{\mathrm{fam}}=3$ from Step~2 and $(b,s)=(3,2)$ from the present step, one recovers
\[
\Oadm
=
N_{\mathrm{fam}}\dim S^+
=
3\cdot 16
=
48
=
(g-1)(b^2+s^2-1)
=
4\cdot 12.
\]

\textbf{Step 4: hypercharge and one-family packet.}
With $(b,s)=(3,2)$ fixed, the even half-spin packet is
\[
S^+=\Lambda^{\mathrm{even}}(E_3\oplus E_2).
\]
Primitive unimodularity on the two-eigenvalue carrier gives
\[
3q_-+2q_+=0.
\]
The unique primitive integer solution with $q_-<0<q_+$ is
\[
(q_-,q_+)=(-2,3).
\]
Hence
\[
X=-2P_3+3P_2,
\qquad
Y=\frac{X}{6}
=
\mathrm{diag}\!\left(-\frac13,-\frac13,-\frac13,\frac12,\frac12\right).
\]

\textbf{Step 5: determinant class and strong CP.}
The determinant-line data are classified by the clutching pair $(c_1(L_2),c_1(L_3))$. On the
admissible branch the compact bosonic index satisfies
\[
\Ind(\mathcal B_{E_2}^+)\ge 0,
\qquad
\Ind(\mathcal B_{E_3}^+)\ge 0.
\]
Existence of a seam-even Higgs doublet forces $\Ind(\mathcal B_{E_2}^+)\ge 1$, while absence of an
unavoidable light color triplet forces $\Ind(\mathcal B_{E_3}^+)=0$. Minimal clutching therefore
gives the unique pair
\[
(c_1(L_2),c_1(L_3))=(1,0).
\]
For the determinant-phase sector, reflection positivity and determinant-line conjugation symmetry
make the admissible free energy even in $\theta_{\mathrm{eff}}$, and the positive-transfer theorem
gives a unique minimum at the origin. Hence
\[
\theta_{\mathrm{eff}}=0.
\]

All discrete blocks are therefore fixed jointly by the same boundary primitive kernel. Thus
\[
\mathfrak D_{\mathrm{disc}}(\Tker^\partial)=\{d^\star_{\mathrm{disc}}\}.
\]
\end{proof}

\begin{corollary}[Derived post-carrier kernel]
\label{cor:derived-post-carrier-kernel}
The tuple
\[
\Tker
:=
(\mathcal A,\mathcal H,D,J,\Gamma,\iC,\Pprim,\Padm,E_3\oplus E_2,Y,[u_\Sigma],\cthree)
\]
is a theorem-level derived object of $\Tker^\partial$. In particular the old post-carrier kernel is
no longer primitive; it is reconstructed from the unique discrete datum $d^\star_{\mathrm{disc}}$
together with the full selector $\Padm=\Pprim P_{\mathrm{sing}}P_\Theta$.
\end{corollary}

\begin{proof}
By \Cref{thm:joint-discrete-admissible-sector}, the carrier split $E_3\oplus E_2$, the charge
operator $Y$, the determinant class, and the family geometry are already fixed by
$\Tker^\partial$. The projector $P_{\mathrm{sing}}$ is then determined by the fixed carrier, and
$P_\Theta$ is determined by the fixed determinant class. Therefore
\[
\Padm=\Pprim P_{\mathrm{sing}}P_\Theta
\]
is fully determined. The displayed tuple is thus derived, not primitive.
\end{proof}

\begin{definition}[Projected modular free energy and stratified master functional]
\label{def:projected-modular-free-energy}
Fix the unique joint discrete datum $d^\star_{\mathrm{disc}}$ of
\Cref{thm:joint-discrete-admissible-sector}, and let $\Padm$ be the derived full selector of
\Cref{cor:derived-post-carrier-kernel}. For
\[
\xi=(\alpha,\chigeo,\delta_{\mathrm{ph}},\rho_{\mathrm{vac}})
\]
on the retained admissible branch, define
\[
\begin{aligned}
\mathcal M_{d^\star_{\mathrm{disc}}}(\xi)
:={}&-\log Z^{\Padm}_{\mathrm{rel}}[\xi]
+\Re\log\det_{\zeta,\Padm}(D_\xi^2+\mu^2)\\
&+\frac12\eta_{\Sigma,\Padm}(D_\xi)
-\log\det(1-U_{6,\xi})
+\Gamma_{\mathrm{vac}}(\rho_\xi),
\end{aligned}
\]
where every term is evaluated on the retained admissible sector $\operatorname{Ran}(\Padm)$.
For a general admissible spectral-bordism state
\[
\mathfrak s=(\mathcal B,\xi),
\qquad
\mathcal B\in\mathbf{SBord}^{\mathrm{adm}},
\]
define the lexicographic barrier
\[
\mathbb B_{\mathrm{lex}}(\mathfrak s)
:=
\begin{cases}
0,& \mathcal B\cong \mathcal B_{\min},\\
+\infty,& \text{otherwise},
\end{cases}
\]
and the stratified master functional
\[
\mathbb M(\mathfrak s)
:=
\mathbb B_{\mathrm{lex}}(\mathfrak s)+\mathcal M_{d^\star_{\mathrm{disc}}}(\xi).
\]
Thus the continuous retained-branch functional $\mathcal M_{d^\star_{\mathrm{disc}}}$ is the
restriction of one stratified functional to the unique minimal admissible stratum.
\end{definition}

\begin{theorem}[Stratified master functional and continuous closure]
\label{thm:continuous-closure-from-master-functional}
On the retained admissible branch, the projected modular free energy factorizes as
\[
\mathcal M_{d^\star_{\mathrm{disc}}}(\xi)
=
\Gamma_{U(1)}(\alpha)
+
\Gamma_{\mathrm{grav}}(\chigeo)
+
\Gamma_{\mathrm{tr}}(\delta_{\mathrm{ph}})
+
\Gamma_{\mathrm{vac}}(\rho_{\mathrm{vac}})
+\mathrm{const}.
\]
Its admissible critical points are therefore exactly the solutions of
\[
\partial_\alpha\Gamma_{U(1)}(\alpha)=0,
\qquad
\partial_{\chigeo}\Gamma_{\mathrm{grav}}(\chigeo)=0,
\qquad
\partial_\delta\Gamma_{\mathrm{tr}}(\delta_{\mathrm{ph}})=0,
\qquad
\delta_\rho\Gamma_{\mathrm{vac}}(\rho_{\mathrm{vac}})=0.
\]
Equivalently, the continuous closure laws are
\[
F_{U(1)}(\alpha)=0,
\qquad
\mathcal R_{\mathrm{grav}}=0,
\qquad
P_{\mathrm{tr}}(\delta_{\mathrm{ph}})=0,
\qquad
\delta_\rho\Gamma_{\mathrm{vac}}(\rho_{\mathrm{vac}})=0.
\]
In particular the continuous closure equations are derived from one primitive variational object
rather than imposed as independent sector primitives. Equivalently, the stationary points of the
stratified master functional $\mathbb M$ are exactly the stationary points of
$\mathcal M_{d^\star_{\mathrm{disc}}}$ on the minimal admissible stratum.
\end{theorem}

\begin{proof}
By construction,
\[
\mathbb M(\mathfrak s)=+\infty
\]
off the unitary-equivalence class of the minimal admissible bordism $\mathcal B_{\min}$. Hence every
stationary point of $\mathbb M$ must lie on that minimal stratum, where the barrier vanishes and
\[
\mathbb M(\mathfrak s)=\mathcal M_{d^\star_{\mathrm{disc}}}(\xi).
\]

Now apply \Cref{cor:derived-post-carrier-kernel}: all continuous terms are evaluated on one and the
same retained sector $\operatorname{Ran}(\Padm)$. On that branch the family, determinant, transport,
and vacuum projectors commute, and the physical Hilbert space splits into orthogonal direct summands
corresponding to the electromagnetic, geometric, transport, and vacuum blocks. For a block-diagonal
operator, the relative partition function factorizes, the zeta determinant factorizes, and the eta
invariant is additive. Therefore the logarithms split $\mathcal M_{d^\star_{\mathrm{disc}}}$ into a
sum of four sector functionals plus a constant.

The $U(1)$ block is exactly the already-defined electromagnetic effective action, so
\[
\partial_\alpha\Gamma_{U(1)}(\alpha)=F_{U(1)}(\alpha).
\]
The geometric block is the reduced spectral Einstein functional, hence
\[
\partial_{\chigeo}\Gamma_{\mathrm{grav}}(\chigeo)=\mathcal R_{\mathrm{grav}}.
\]
The transport block is the logarithmic resolvent determinant of the exact $C_6$ transfer
denominator, so
\[
\partial_\delta\Gamma_{\mathrm{tr}}(\delta_{\mathrm{ph}})
=
P_{\mathrm{tr}}(\delta_{\mathrm{ph}}).
\]
Finally the vacuum block is by definition the admissible vacuum effective action, so
\[
\delta_\rho\Gamma_{\mathrm{vac}}(\rho_{\mathrm{vac}})=0
\]
is exactly its Euler--Lagrange equation. Thus the four continuous closure equations are the
Euler--Lagrange equations of one projected modular free energy.
\end{proof}

\begin{definition}[Sector solution sets]
\label{def:intrinsic-closure-map}
\label{def:core-closure-functional}
\label{def:sector-solution-sets}
Define
\[
Z_{\mathrm{disc}}:=\mathfrak D_{\mathrm{disc}}(\Tker^\partial),
\qquad
Z_{U(1)}:=\{\alpha:\ F_{U(1)}(\alpha)=0\},
\qquad
Z_{\mathrm{grav}}:=\{\chigeo:\ \mathcal R_{\mathrm{grav}}(\chigeo)=0\},
\]
\[
Z_{\mathrm{tr}}:=\{\delta_{\mathrm{ph}}:\ P_{\mathrm{tr}}(\delta_{\mathrm{ph}})=0\},
\qquad
Z_{\mathrm{vac}}:=\{\rho_{\mathrm{vac}}:\ \delta_\rho\Gamma_{\mathrm{vac}}(\rho_{\mathrm{vac}})=0\}.
\]
Let
\[
\mathcal I_0
:=
(\Padm,\iC,[u_\Sigma],E_3\oplus E_2,Y,\cthree)
\]
be the upstream invariant tuple fixed before the continuous closure step.
\end{definition}

\begin{theorem}[Intrinsic branch as a singleton fiber product]
\label{thm:intrinsic-closed-branch}
\label{thm:observation-closure}
\label{thm:singleton-fiber-product-branch}
For every admissible boundary primitive kernel $\Tker^\partial$, the intrinsic closed branch is the
singleton fiber product
\[
\mathcal Z_{\mathrm{cl}}
\cong
Z_{\mathrm{disc}}
\times_{\mathcal I_0}
Z_{U(1)}
\times_{\mathcal I_0}
Z_{\mathrm{grav}}
\times_{\mathcal I_0}
Z_{\mathrm{tr}}
\times_{\mathcal I_0}
Z_{\mathrm{vac}},
\]
and therefore
\[
\mathcal Z_{\mathrm{cl}}=\{x_{\mathrm{cl}}^\star\}.
\]
No Jacobian on mixed discrete and continuous variables is needed.
\end{theorem}

\begin{proof}
By \Cref{thm:joint-discrete-admissible-sector},
\[
Z_{\mathrm{disc}}=\{d^\star_{\mathrm{disc}}\}.
\]
By the exact electromagnetic theorem,
\[
Z_{U(1)}=\{\alpha_\star\}.
\]
By coercivity and strict convexity of the reduced Einstein functional,
\[
Z_{\mathrm{grav}}=\{\chigeo^\star\}.
\]
By the algebraic transport theorem together with the word-length lock,
\[
Z_{\mathrm{tr}}=\{\delta_{\mathrm{ph}}^\star\}.
\]
By the positive-transfer, thermodynamic-limit, and gap-stability theorems,
\[
Z_{\mathrm{vac}}=\{\rho_{\mathrm{vac}}^\star\}.
\]

Each singleton carries its compatibility map to the same upstream invariant tuple $\mathcal I_0$.
That common tuple is already fixed earlier: the discrete theorem fixes $E_3\oplus E_2$ and $Y$,
the derived selector theorem fixes $\Padm$, and the seam theorem fixes $[u_\Sigma]$ and $\cthree$.
Hence all compatibility maps land at the same point of $\mathcal I_0$. The fiber product of
singleton sets over one common point is again a singleton. Therefore
\[
\mathcal Z_{\mathrm{cl}}=\{x_{\mathrm{cl}}^\star\}.
\]
\end{proof}

\begin{theorem}[Internal analytic closure as a corollary of primitive completion]
\label{thm:internal-analytic-closure}
\label{thm:internal-analytic-closure-hard}
For every admissible one-sided boundary datum, the analytic closure blocks needed for rigid
reconstruction hold internally on the unique intrinsic branch determined by
\Cref{thm:joint-discrete-admissible-sector,thm:continuous-closure-from-master-functional,thm:singleton-fiber-product-branch}.
\end{theorem}

\begin{proof}
The Calder\'on and APS block is already fixed upstream by the primitive boundary theorem chain.
The family, determinant, and effective strong-angle blocks are fixed by the joint discrete theorem.
The electromagnetic, geometric, transport, and vacuum blocks are fixed by the continuous master
functional together with the already-proved positivity, gap, and thermodynamic-limit theorems.
Because the intrinsic branch is a singleton by \Cref{thm:singleton-fiber-product-branch}, all
analytic closure blocks hold on one and the same branch, internally and without any extra closure
package.
\end{proof}

\begin{theorem}[Canonical rigid reconstruction and essential injectivity on the rigid image]
\label{thm:rigid-reconstruction-core}
\label{thm:canonical-rigid-reconstruction-essential-injective}
Let $\Tbound$ be an admissible one-sided boundary datum and let $\Tker^\partial$ be its boundary
primitive kernel. Then the rigid tuple
\[
\mathfrak C(\Tker^\partial)
:=
\bigl(
E_3\oplus E_2,\,
Y,\,
S^+,\,
X_f^\circ,\,
[\nabla_F^\star],\,
\Phi^\star,\,
\chigeo^\star,\,
\delta_{\mathrm{ph}}^\star,\,
\rho_{\mathrm{vac}}^\star,\,
\theta_{\mathrm{eff}}=0,\,
U_\Sigma^\star,\,
\Lambda_{\mathrm{IR}}^\star
\bigr)
\]
is uniquely determined. In particular the reconstruction functor is essentially injective on
objects on its rigid image, and every two rigid realizations coming from the same boundary datum are
uniquely unitarily isomorphic.
\end{theorem}

\begin{proof}
The joint discrete theorem fixes $E_3\oplus E_2$, $Y$, $S^+$, $X_f^\circ$, $[\nabla_F^\star]$, the
determinant class, and $\theta_{\mathrm{eff}}=0$. The seam theorem fixes $[u_\Sigma]$ and
$\cthree$. The continuous master-functional theorem fixes $\alpha_\star$, $\chigeo^\star$,
$\delta_{\mathrm{ph}}^\star$, and $\rho_{\mathrm{vac}}^\star$. The downstream definitions of
$\Phi^\star$, $U_\Sigma^\star$, and $\Lambda_{\mathrm{IR}}^\star$ are then evaluated on that unique
branch. Hence every displayed entry of $\mathfrak C(\Tker^\partial)$ is uniquely determined.

If two rigid realizations arise from the same boundary datum, then all displayed entries agree. The
carrier, seam, family, determinant, transport, and vacuum blocks each admit only the unique
unitary intertwiner preserving the already-fixed theorem-level structures, and these intertwiners
are compatible because they preserve the same upstream invariant tuple
\[
(\Padm,\iC,[u_\Sigma],E_3\oplus E_2,Y,\cthree).
\]
Their direct sum or tensor product therefore yields one global unitary intertwiner. Thus the rigid
image is essentially injective on objects, with unique unitary isomorphism on that image.
\end{proof}
